%% Beamer Template of Sabine Hauert and R.~Eddie Wilson
%% Available on https://www.overleaf.com/latex/templates/eddies-minimum-beamer-style/qhcyjdqkbbgx#.V-vfg1SLS1I

%%%%%%%%%%%%%%%%%%%%%%%%%%%%%%%%%%%%%%%%%%%%%%
% Head matter - can we try to be consistent on
% included packages

\documentclass{beamer}
\mode<presentation>
{\usetheme{default}
 \usecolortheme{default}
 \usefonttheme{default}
 \setbeamertemplate{navigation symbols}{}
 \setbeamertemplate{caption}[numbered]} 


% Block colors
\setbeamercolor{block body}{bg=blue!10,fg=black}
\setbeamercolor{block title}{bg=blue!30,fg=black}

% Pour que tikzcd fonctionne
%\[\begin{tikzcd}[ampersand replacement=\&, column sep=small]\end{tikzcd}\]


\usepackage[frenchb]{babel}
\usepackage{amsfonts}
\usepackage{amsmath}
\usepackage{amssymb}
%\usepackage[T1]{fontenc}
\usepackage[utf8]{inputenc}
\usepackage{amsthm}
\usepackage{graphicx}
\usepackage{tikz}
\usepackage{tikz-cd}
\usepackage{hyperref}
\usepackage{amssymb}
\usepackage{geometry}

\hypersetup{                    % parametrage des hyperliens
    colorlinks=true,                % colorise les liens
    breaklinks=true,                % permet les retours à la ligne pour les liens trop longs
    urlcolor= blue,                 % couleur des hyperliens
    linkcolor= blue,                % couleur des liens internes aux documents (index, figures, tableaux, equations,...)
    citecolor= cyan               % couleur des liens vers les references bibliographiques
    }

\theoremstyle{definition}
\newtheorem{definition}{Definition}
\newtheorem{thm}{Theorem}
\newtheorem{ex}{Exercice}
\newtheorem{lem}{Lemma}
\newtheorem*{dem}{Proof}
\newtheorem{prop}{Proposition}
\newtheorem{cor}{Corollary}
\newtheorem{conj}{Conjecture}
\newtheorem{Res}{Result}
\newtheorem{Expl}{Example}
\newtheorem{rk}{Remark}

\newcommand{\N}{\mathbb N}
\newcommand{\Z}{\mathbb Z}
\newcommand{\R}{\mathbb R}
\newcommand{\C}{\mathbb C}
\newcommand{\Hil}{\mathcal H}
\newcommand{\Mn}{\mathcal M _n (\mathbb C)}
\newcommand{\K}{\mathbb K}
\newcommand{\B}{\mathbb B}
\newcommand{\Cat}{\mathbb B / \mathbb K}
\newcommand{\G}{\mathcal G }

\setlength\parindent{0pt}


%%%%%%%%%%%%%%%%%%%%%%%%%%%%%%%%%%%%%%%%%%%%%%
% Formatting for title page
\title[First Steps with SCRATCH]{$K$-théorie contrôlée pour les groupoïdes\\ et applications}
\author{Clément Dell'Aiera}
\institute{IECL}
\date{12 Juillet 2017}
%%%%%%%%%%%%%%%%%%%%%%%%%%%%%%%%%%%%%%%%%%%%%%

\begin{document}

%%% TITRE
\begin{frame}
  \titlepage
\begin{center}\includegraphics[width=5cm]{IECL.png}\end{center}
\end{frame}

\begin{frame}
  \tableofcontents
\end{frame}

\section{Introduction}
\begin{frame}{Introduction}
En physique, les modèles prévoient les valeurs de quantités mesurables, appelées \textbf{observables}. Ce sont des fonctions
\[f:\Omega \rightarrow \R\]
où $\Omega$ est l'espace de configuration du système. Ces observables commutent en physique classique.\\
\vspace{0.3 cm} 
Pour donner un modèle de l'atome d'hydrogène en accord avec l'expérience, Heisenberg propose en 1925 de remplacer les observables classiques par des "q-nombres", qui ne commutent pas.
\[pq -qp = i\hbar 1\]
\end{frame}

\begin{frame}{Introduction}
Les "nombres" utilisés en physique quantique forment une algèbre d'opérateurs. Ces algèbres sont encodées d'un point de vue abstrait par la notion de $C^*$-algèbre. 
\vspace{0.3 cm}
\begin{definitionfr}
Une $C^*$-algèbre est une $\C$-algèbre $A$ munie d'un antihomomorphisme involutif $*$ et d'une norme multiplicative $||.||$ telle que :
\begin{itemize}
\item[$\bullet$] $(A,||.||)$ est une algèbre de Banach,
\item[$\bullet$] $||a^*a||^2 = ||a||^2$ pour tout $a\in A$.
\end{itemize}
\end{definitionfr}

\begin{itemize}
\item[$\bullet$] l'algèbre des opérateurs bornés sur un espace de Hilbert $\mathcal L(H)$, 
\item[$\bullet$] l'algèbre des fonctions continues sur un espace localement compact $C_0(X)$.
\end{itemize}

\end{frame}

\begin{frame}{Introduction}
Il s'avère que toute $C^*$-algèbre commutative est équivalente à un espace localement compact.
\begin{thmfr}[Dualité de Gelfand]
Soit $A$ une $C^*$-algèbre commutative. Alors il existe un espace localement compact $X$ et un $*$-isomorphisme 
\[A\cong C_0(X).\]
\end{thmfr}

Alain Connes développe à partir des années 80 la géométrie non-commutative basée sur le principe suivant
\begin{block}{Géométrie Non-Commutative}
Une $C^*$-algèbre représente un "espace non-commutatif". 
\end{block}

\end{frame}

\begin{frame}{Introduction}
L'idée des physiciens (remplacer l'algèbre des coordonnées par une algèbre d'opérateurs) peut être utilisée pour étudier des espaces pathologiques.
\vspace{0.3 cm}
\begin{block}{Philosophie}
Associer à une situation géométrique une $C^*$-algèbre.
\end{block}
\vspace{0.3 cm}
Quels espaces pathologiques ? Principalement, des espaces d'orbites, associés à :
\begin{itemize}
\item[$\bullet$] des actions de groupes,
\item[$\bullet$] des feuilletages,
\item[$\bullet$] des sytèmes dynamiques.
\end{itemize} 
\end{frame}

\begin{frame}{Introduction}
\textbf{Exemple :} Le groupe $\Z$ agit sur le cercle $\mathbb S^1$ par rotation d'angle $2\pi\theta$.
\vspace{0.3 cm}
\[\begin{tikzpicture}
\draw[fill = gray!20] (0,0) circle (2) ;
\draw [>=stealth, ->, magenta] (2.2,0)  arc (0:60:2.2); %--(1 , 1.72);
\draw[magenta]  (2.3,1.5) node {$\theta$};
\draw  (2,0) node {$\bullet$};
\draw  (1,1.72) node {$\bullet$};
\end{tikzpicture}\]
\end{frame}

\begin{frame}{Introduction}
Deux cas :
\vspace{0.3 cm}
\begin{itemize}
\item[$\bullet$] $\theta\in \mathbb Q$ : les orbites sont périodiques, et le quotient est 
\[\mathbb S^1 / \Z \cong \mathbb S^1,\] 
\item[$\bullet$] $\theta \notin \mathbb Q$ : toute orbite est dense, et  $\mathbb S^1 / \Z$ est un espace hautement pathologique. Rien ne le distingue d'un point.
 \[C_0(\mathbb S^1 / \Z) \cong \C.\]
\end{itemize}
\vspace{0.3 cm}
Une solution : construire un groupoïde qui encode la situation géométrique. On sait associer des $C^*$-algèbres aux groupoïdes.\\
\end{frame}

\begin{frame}{Introduction}

\vspace{0.3 cm}
\begin{definitionfr}[Groupoïde]
Un groupoïde est la donnée de deux espaces topologiques, l'espace des flèches $G$ et celui des unités $G^{(0)}$, ainsi que de :
\begin{itemize}
\item[$\bullet$] un plongement topologique $e: G^{(0)} \rightarrow G; x\mapsto e_x$, 
\item[$\bullet$] des applications continues surjectives source et but $s,r : G\rightrightarrows G^{(0)}$ telles que $r\circ e =s\circ e = id_{G^{(0)}}$,
\item[$\bullet$] une multiplication continue \[G^{(2)}=G\times_{s,r} G\rightarrow G; (g,g') \mapsto gg'\] 
telle que $g(g'g'') = (gg')g''$ et $ge_{s(g)} = e_{r(g)}g= g$,
\item[$\bullet$] une involution continue $G\rightarrow G ; g\mapsto g^{-1}$ telle que $gg^{-1} = e_{r(g)}$, $g^{-1}g = e_{s(g)}$ et $s(g^{-1})=r(g)$.	
\end{itemize}
\end{definitionfr}
\end{frame}

\begin{frame}{Introduction}
\[\begin{tikzpicture}
\draw  (0,1) node {$g$};
\draw  (0,-1) node {$g^{-1}$};
\draw [>=stealth, ->, thick] (-1,0.2) to[bend left] (1,0.2);
\draw [>=stealth, ->,thick] (1,-0.2) to[bend left] (-1,-0.2);
\draw  (-1,0) node {$\bullet$};
\draw  (-1.5,0) node {$s(g)$};
\draw  (1,0) node {$\bullet$};
\draw  (1.5,0) node {$r(g)$};
\end{tikzpicture}\]
\end{frame}

\begin{frame}{Introduction}
Des exemples de groupoïdes :
\vspace{0.3 cm}
\begin{itemize}
\item[$\bullet$] un espace topologique $X\rightrightarrows X$,
\item[$\bullet$] un groupe topologique $G\rightrightarrows *$,
\item[$\bullet$] le groupoïde des paires $X\times X\rightrightarrows X$,
\[\begin{tikzpicture}
\draw  (-2,1) node {$(x,y)$};
\draw  (-2,-1) node {$(y,x)$};

\draw  (1,1) node {$(x,y)$};
\draw  (2,1) node {$(y,z)$};
\draw  (1.5,-1) node {$(y,z)$};

\draw [>=stealth, ->, thick] (-2.5,0.2) to[bend left] (-1.5,0.2);
\draw [>=stealth, ->,thick] (-1.5,-0.2) to[bend left] (-2.5,-0.2);

\draw [>=stealth, ->, thick] (0.5,0.2) to[bend left] (1.45,0.2);
\draw [>=stealth, ->, thick] (1.55,0.2) to[bend left] (2.5,0.2);
\draw [>=stealth, ->, thick] (0.5,-0.2) to[bend right] (2.5,-0.2);


\draw  (-2.5,0) node {$\bullet$};
\draw  (-1.5,0) node {$\bullet$};
\draw  (0.5,0) node {$\bullet$};
\draw  (1.5,0) node {$\bullet$};
\draw  (2.5,0) node {$\bullet$};

\draw  (-3,0) node {$x$};
\draw  (-1,0) node {$y$};
\draw  (0,0) node {$x$};
\draw  (1.8,-0.2) node {$y$};
\draw  (3,0) node {$z$};
\end{tikzpicture}\]
\item[$\bullet$] le dernier exemple se généralise aux relations d'équivalence : $R\subseteq X\times X$.
\end{itemize}
\end{frame}

\begin{frame}{Introduction}
$\bullet$ Soit $\Gamma$ un groupe dénombrable discret agissant sur $X$ par homéomorphismes. Le groupoïde produit croisé noté $X\rtimes \Gamma \rightrightarrows X$.
\[\begin{tikzpicture}
\draw  (-2,1) node {$(x,g)$};
\draw  (-2,-1) node {$(gx,g^{-1})$};

\draw  (0.8,1) node {$(x,g)$};
\draw  (2.2,1) node {$(gx,g')$};
\draw  (1.5,-1) node {$(gx,g'gx)$};

\draw [>=stealth, ->, thick] (-2.5,0.2) to[bend left] (-1.5,0.2);
\draw [>=stealth, ->,thick] (-1.5,-0.2) to[bend left] (-2.5,-0.2);

\draw [>=stealth, ->, thick] (0.5,0.2) to[bend left] (1.45,0.2);
\draw [>=stealth, ->, thick] (1.55,0.2) to[bend left] (2.5,0.2);
\draw [>=stealth, ->, thick] (0.5,-0.2) to[bend right] (2.5,-0.2);


\draw  (-2.5,0) node {$\bullet$};
\draw  (-1.5,0) node {$\bullet$};
\draw  (0.5,0) node {$\bullet$};
\draw  (1.5,0) node {$\bullet$};
\draw  (2.5,0) node {$\bullet$};

\draw  (-3,0) node {$x$};
\draw  (-1,0) node {$gx$};
\draw  (0,0) node {$x$};
\draw  (1.8,-0.2) node {$gx$};
\draw  (3,0) node {$g'gx$};
\end{tikzpicture}\]
\end{frame}

\begin{frame}{Introduction}
Quelques $C^*$-algèbres associées :
\vspace{0.3 cm}
\begin{itemize}
\item[$\bullet$] $C_r^*(X)\cong C_0(X)$,
\vspace{0.3 cm}
\item[$\bullet$] dans le cas d'un groupe, on retrouve la $C^*$-algèbre de groupe,
\vspace{0.3 cm}
\item[$\bullet$] si $X$ est fini, $C^*_r(X\times X) \cong \mathcal L (l^2(X)) \cong \mathfrak M_{|X|}(\C)$,
\vspace{0.3 cm}
\item[$\bullet$] $C^*_r(X\rtimes \Gamma) \cong C_0(X)\rtimes_r \Gamma$.
\end{itemize}
\end{frame}

\begin{frame}{Introduction}
Dans le cas de la topologie classique, un des objets à déterminer est ce qu'on appelle l'homologie et la cohomologie d'un espace.\\
\vspace{0.3 cm}
Dans le cadre de la géométrie non-commutative, on voudra déterminer la $K$-théorie des certaines $C^*$-algèbres, qui est l'analogue des théories homologiques classiques. 
\vspace{0.3 cm}
\begin{block}{Objectif de la thèse}
Déterminer les groupes de $K$-théorie de certaines familles de $C^*$-algèbres.
\end{block}
\end{frame}

\begin{frame}{Introduction}
Pour cela, voici notre stratégie :\\
\vspace{0.3 cm}
\begin{itemize}
\item[$\bullet$] détecter la filtration naturelles des produits croisés de $C^*$-algèbres par des groupoîdes étales et des groupes quantiques discrets,
\vspace{0.3 cm}
\item[$\bullet$] approximer la $K$-théorie par la $K$-théorie contrôlée,
\vspace{0.3 cm}
\item[$\bullet$] dans le cas des $C^*$-algèbres de Roe et des produits croisés de groupoïdes étales, construire des applications d'assemblage,
\vspace{0.3 cm}
\item[$\bullet$] prouver une version contrôlée de la conjecture de Baum-Connes.
\end{itemize}
\end{frame}


\section{Motivations}
\begin{frame}
  \tableofcontents[currentsection]
\end{frame}
\begin{frame}{Motivations}

Soit $X$ un espace métrique dénombrable discret à géométrie bornée, i.e. $\sup_{x\in X} |B(x,R)|<\infty$ pour tout $R>0$. On note $H_X = l^2(X)\otimes H$ et 

\[C_R[X] = \{T\in \mathcal L(H_X) \text{ t.q. } T_{xy} \in \mathfrak K(H) \text{ et } prop(T) < R \}\]

\begin{definitionfr}[Algèbre de Roe]
\[C^*(X) = \overline{\cup_{R>0} C_R[X]} \subseteq \mathcal L(H_X).\]
\end{definitionfr}
\end{frame}

\begin{frame}{Motivations}
\begin{conj}[Baum-Connes coarse]
Soit $X$ un espace métrique dénombrable discret à géométrie bornée. Alors l'application d'assemblage
\[\mu_X : \varinjlim KK_*(C_0(P_d(X),\C) \rightarrow K_*(C^*(X))\]
est un isomorphisme.
\end{conj}

Soit $\Gamma$ un groupe finiment engendré par une partie symmétrique $S$. \\
\vspace{0.3 cm}
On note $|\Gamma|$ l'espace métrique obtenu en munissant $\Gamma$ de la métrique de la longueur des mots. \\
\vspace{0.3 cm}
La conjecture de Baum-Connes coarse pour $|\Gamma |$ implique la conjecture de Novikov sur les hautes signatures associées à $\Gamma$. 
\end{frame}

\begin{frame}{Motivations}

\begin{thmfr}[Yu, 2010 \cite{Yu1}]
Soit $X$ de dimension asymptotique finie. Alors l'application d'assemblage
\[\mu_X : \varinjlim KK(C_0(P_d(X),\C) \rightarrow K(C^*(X))\]
est un isomorphisme.
\end{thmfr}

Dans la preuve : des groupes notés $QP_{\delta, r ,s, k}(X)$ et $QU_{\delta, r ,s, k}(X)$ qui approximent $K_0(C^*(X))$ et $K_1(C^*(X))$.\\

\end{frame}

\begin{frame}{Motivations}

Un cadre plus général : la $K$-théorie quantitative (Oyono-Oyono \& Yu, 2015 \cite{OY2}) pour les $C^*$-algèbres filtrées.\\
\vspace{0.3 cm}
Une preuve plus générale :

\begin{thmfr}[Oyono-Oyono Yu, 2017]
Soit $X$ de décomposition géométrique finie. Alors l'application d'assemblage
\[\mu_X : \varinjlim KK(C_0(P_d(X),\C) \rightarrow K(C^*(X))\]
est un isomorphisme.
\end{thmfr}

\end{frame}


\section{$K$-théorie contrôlée}
\begin{frame}
  \tableofcontents[currentsection]
\end{frame}

\begin{frame}{$K$-théorie contrôllée}

\begin{definitionfr}
Une structure coarse $\mathcal E$ est un semi-groupe abélien qui est aussi un treillis. \\
\end{definitionfr}
\vspace{0.3 cm}
Rappelons qu'un treillis est un ensemble partiellement ordonné tel que toute paire $(E,E')$ admette un supremum $E\vee E'$ et un infimum $E\wedge E'$.\\
\vspace{0.3 cm}
On note $(E,E')\mapsto E\circ E'$ la loi de composition de $\mathcal E$.
\end{frame}

\begin{frame}{$K$-théorie contrôllée}
\begin{definitionfr}
Une $C^*$-algèbre $A$ est $\mathcal E$-filtrée s'il existe une famille de sous-espaces vectoriels fermés auto-adjoints $\{A_E\}_{E\in\mathcal E}$ de $A$ telle que :
\begin{itemize}
\item[$\bullet$] si $E\leq E'$, alors $A_E\subseteq A_{E'}$,
\item[$\bullet$] pour tout $E,E'\in\mathcal E$, $A_E.A_{E'}\subseteq A_{E\circ E'}$,
\item[$\bullet$] $A$ est l'adhérence de l'union des sous-espaces $A_E$, i.e. $\overline{\cup_{E\in\mathcal E}A_E} = A$.
\item[$\bullet$] si $A$ est unitale, on a de plus $1\in A_E,\forall E\in\mathcal E$.
\end{itemize}
\end{definitionfr}

Si $\mathcal E = \R_+^*$ muni de la somme, on retrouve la définition de Oyono-Oyono et Yu. 
\end{frame}

\begin{frame}{$K$-théorie contrôllée}

Plus de $C^*$-algèbres peuvent être vues commes des $C^*$-algèbres filtrées.

\begin{exple}
Les $C^*$-algèbres de Roe sont filtrées.
\end{exple}

\begin{exple}
Les produits croisés de $C^*$-algèbres par des actions de groupoïdes étales sont filtrées.
\end{exple}

\begin{exple}
Les produits croisés de $C^*$-algèbres par des actions de groupes quantiques discrets sont filtrées.
\end{exple}

\end{frame}

\begin{frame}{$K$-théorie contrôllée}
Soit $\mathcal E$ une structure coarse et $A$ une $C^*$-algèbre $\mathcal E$ filtrée. On définit les $\varepsilon$-$E$-unitaires
\[U^{\varepsilon, E}(A)= \{u\in A_E \text{ s.t. } ||u^*u-1||<\varepsilon\text{ and }||uu^*-1||<\varepsilon \}\]
et les $\varepsilon$-$E$-projections 
\[P^{\varepsilon, E}(A)= \{p\in A_E \text{ s.t. } p=p^*\text{ and }||p^2-p||<\varepsilon \}.\]
\end{frame}

\begin{frame}{$K$-théorie contrôllée}
Comme en $K$-théorie, 
\begin{itemize}
\item[$\bullet$] $P_\infty^{\varepsilon, E}(A)$ est la limite inductive algébrique des $P_n^{\varepsilon, E}(A)$ par rapports aux inclusions
\[\left\{\begin{array}{rcl}
	P^{\varepsilon,E}_n(A) 		& \rightarrow	& P^{\varepsilon,E}_{n+1}(A)\\ 
	p 		& \mapsto 	& \begin{pmatrix}p& 0 \\ 0&0 \end{pmatrix}
\end{array}\right.\]
\item[$\bullet$] $U_\infty^{\varepsilon, E}(A)$ est la limite inductive algébrique des $U_n^{\varepsilon, E}(A)$ par rapports à
\[\left\{\begin{array}{rcl}
	U^{\varepsilon,E}_n(A) 		& \rightarrow	& U^{\varepsilon,E}_{n+1}(A)\\ 
	u 		& \mapsto 	& \begin{pmatrix}u & 0 \\ 0& 1 \end{pmatrix}
\end{array}\right. .\]
\end{itemize}
\end{frame}

\begin{frame}{$K$-théorie contrôllée}
On munit $P_\infty^{\varepsilon, E}(A)\times \N$ et $U_\infty^{\varepsilon, E}(A)$ des relations d'équivalence suivantes:
\begin{itemize}
\item[$\bullet$] $(p,l) \sim (q,l')$ s'il existe une homotopie de quasi-projections $h\in P^{\varepsilon, E}_\infty(A[0,1])$ et un entier $k$ tel que 
\[h(0)=\begin{pmatrix} p & 0 \\ 0 & 1_{k+l'} \end{pmatrix} \text{ and }
h(1)=\begin{pmatrix} q & 0 \\ 0 & 1_{k+l} \end{pmatrix}\]
\item[$\bullet$] $u \sim v$ s'il existe une homotopie de quasi-unitaires $h\in U^{3\varepsilon, E\circ E}_\infty(A[0,1])$ tel que $h(0)= u \text{ and }h(1)=v$.\\
\end{itemize}
\end{frame}

\begin{frame}{$K$-théorie contrôllée}
Si $A$ est unitale,
\begin{itemize}
\item[$\bullet$] $K_0^{\varepsilon,E}(A) = P^{\varepsilon, E}_\infty(A)\times \N / \sim_{\varepsilon,E}$ 
\item[$\bullet$] $K_1^{\varepsilon,E}(A) = U^{\varepsilon, E}_\infty(A) / \sim_{\varepsilon,E}$.
\end{itemize}

Si $A$ n'est pas unitale,  
\[K_0^{\varepsilon,E}(A) = \{[p,l]_{\varepsilon,E} : p\in P^{\varepsilon,E}_\infty (\tilde A), l\in \N \text{ s.t. rank}(\kappa_0(\rho_A(p)))=l \}\]
et $K_1^{\varepsilon,E}(A)$ est défini par $U_\infty^{\varepsilon,E}(\tilde A)/ \sim_{\varepsilon,E}$.\\

\begin{definitionfr}
La $K$-théorie contrôllée d'une $C^*$-algèbre filtrée $(A,\mathcal E)$ est la famille de groupes abéliens 
\[\hat K_0(A) = (K_0^{\varepsilon,E}(A))_{\varepsilon\in (0,\frac{1}{4}),E\in\mathcal E} \text{ et } \hat K_1(A) = (K_1^{\varepsilon,E}(A))_{\varepsilon\in (0,\frac{1}{4}),E\in\mathcal E}.\]
\end{definitionfr}
\end{frame}















































%\section{AssemblyMaps}
%\section{Assembly Maps}

In this section, all groupoids $G\rightrightarrows G^{(0)}$ are étale. 

\subsection{Proper groupoids and proper actions}

\begin{definition}
A topological groupoid $G\rightrightarrows G^{(0)}$ is proper if the map $G\times_s G^{(0)}\rightarrow G^{(0)}\times G^{(0)}$ defined by $(g,z)\mapsto (z,gz)$ is a proper map.\\
A $G$-space $Z$ is said to be proper if the crossed-product groupoid $Z\rtimes G$ is proper. If the space of orbits $Z/G$ is compact, $Z$ is called $G$-compact.\\
\end{definition}

\begin{definition}
A $G$-space $\underline E G$ is called universal if, given any proper $G$-compact $G$-space $Z$, there exists a continuous $G$-equivariant map $Z\rightarrow \underline E G$.
\end{definition}

We now construct a model for universal spaces of $G$, based on J.-L.Tu's work \cite{TuBC2}. For any compact subset $K\subset G$, define $P_K(G)$ to be the space of probability measures $\nu $ with support contained in one and only one fiber $G^x$ for some $x\in G^{(0)}$, and such that if $g,g'\in \text{supp }(\nu)$, then $g'g^{-1}\in K$. We endow $P_K(G)$ with the weak-$*$ topology.\\

The action of $G$ is defined by translation. The momentum map $P_K(G)\rightarrow G^{(0)}$ is just the map associating to $\nu$ the only $x$ such that $\text{supp }(\nu) \subset G^x$. As the fibers are discrete, any $\nu\in P_K(G)$ can be represented as as sum $\nu = \sum_{g\in G^x} \lambda_g(\nu)\delta_g$, where $\delta_g$ is the Dirac measure at $g\in G $. The continuous functions $\lambda_g$ are called coordinate functions and satisfy $\sum_{g\in G^x} \lambda_g(\nu) =1$ for every $x\in G^{(0)}$. The action of $G$ is given by $g\lambda_h = \lambda_{g^{-1}h}$.\\

\begin{lem}[Tu,\cite{TuBC2}]
The action of $G$ on $P_K(G)$ is proper and cocompact.
\end{lem}  

\begin{lem}[Tu,\cite{TuBC2}]
Let $Z$ be a proper $G$-compact $G$-space. Then there exists a compact subset $K\subset G$ and a $G$-equivariant continuous map $Z\rightarrow P_K(G)$.
\end{lem}  

\subsection{Equivariant $K$-homology}

We will use the equivariant $KK$-theory developped by Le Gall in his thesis \cite{LeGall}, which is an extension of the usual equivariant $KK$-theory of Kasparov. Recall that, if $A$ and $B$ are two $G$-algebras, elements of $KK^G(A,B)$ are homotopy classes of triple $(E,\pi,T)$ where :\\

\begin{itemize}
\item[$\bullet$] $E$ is a $G$-module,
\item[$\bullet$] $\pi : A\rightarrow \mathcal L_B(E)$ is a $*$-homomorphism,
\item[$\bullet$] $T\in \mathcal L_B(E)$ is an adjoinable operator such that the triple satisfies the condition of $K$-cycle : $\pi(a)(T^2-T),\pi(a)(T^*-T)[\pi(a),T]$, and $\pi(a)(T-g.T)$ are compact operators if $E$, for all $a\in A, g\in G$.\\
\end{itemize}

\begin{definition}
Let $Y$ be a proper $G$-space and $B$ a $G$-algebra. Then the analytic $K$-homology of $Y$ with coefficients in $B$ is defined by 
\[RK^G(Y,B) = \varinjlim_{Z\subset Y} KK^G(C_0(Z),B), \]
the inductive limit being taken on proper $G$-compact $G$-subspaces $Z$ of $Y$.
\end{definition}

The previous lemmas assures that this inductive limit can be somehow restricted when it comes to the analytic $>K$-homology of any universal space $\underline E G$ :
\[RK^G(\underline E G,B) = \varinjlim_{K\subset G \text{ compact }} KK^G(C_0(P_K(G)),B).\]

\subsection{Descent functor}

There exists a natural transformation $KK^G(A,B)\rightarrow KK(A\rtimes_r G, B\rtimes_r G)$ respecting the Kasparov product, which is called the descent functor. The same statement remains true for maximal crossed products. We recall its construction, which was first stated in \cite{LeGall}.\\
 
Let $A$ and $B$ be two $G$-algebras and $(E,\pi,T)\in E^G(A,B)$ be a $K$-cycle. Define :
\begin{itemize}
\item[$\bullet$] $E_G$ to be the completion of $C_c(G,E)$ with respect to the norm $||f||= \sup_{x\in G^{(0)}} \sum_{G^x} |f(g)|^2$,
\item[$\bullet$] $\pi_G$ to be the image of $\pi:A \rightarrow \mathcal L_B(E)$ under the reduced crossed-product functor,
\item[$\bullet$] $T_G$ to be the image of $T\in\mathcal L(E)$ under the reduced crossed-product functor.
\end{itemize}
Then $(E_G,\pi_G,T_G)\in E(A\rtimes_r G,B\rtimes_r G)$ and the map $(E,\pi,T)\mapsto (E_G,\pi_G,T_G) $ induces a homorphism of abelian groups
\[j_G : KK^G(A,B)\rightarrow KK(A\rtimes_r G,B\rtimes_r G)\]
satisfying $j_G(z\otimes_D z')=j_G(z)\otimes_{D\rtimes_r G} j_G(z')$ for any $z\in KK^G(A,D), z'\in KK^G(D,B)$.\\ 

\subsection{The assembly map}

If $Z$ is a proper $G$-compact $G$-space, $Z\rtimes G$ is a proper groupoid by definition, hence there exists a cutoff function $c : Z\rightarrow [0,1]$ such that $\int_{G^{p(z)}} c(zg)d\lambda^x(g) = 1$. The function $g\mapsto c(r(g))^{\frac{1}{2}}c(s(g))^{\frac{1}{2}} $ defines a projection in $C_0(Z)\rtimes_r G$ which we denote by $\mathcal L_Z$. If $Z=P_K(G)$, then $\mathcal L_Z = \mathcal L_K$.

\begin{definition}
The assembly map for $G$ with coefficients in $B$ is defined as the inductive limit of the maps $\mu_{G,B}^{(Z)} : KK^G(C_0(Z),B)\rightarrow K(B\rtimes_r B)$ given by
\[\mu_{G,B}^{(Z)} (z)=[\mathcal L_Z]\otimes_{C_0(Z)\rtimes G} j_G(z),\]
that is $\mu_{G,B} = \varinjlim \mu_{G,B}^{(Z)}$ (one has to check that theses maps respects the inductive systems, which they do).
\end{definition}

%Another definition : with $Ad$...

\subsection{The Baum-Connes conjecture}

\begin{conj}[Baum-Connes conjecture]
Let $G$ be an étale groupoid, and $A$ a $G$-algebra.\\
The Baum-Connes conjecture for $G$ with coefficients in $A$ is the following claim : \\
\center{$\textbf{BC}_{G,A}$ : $\mu_{G,A}$ is an isomorphism of abelian groups.}\\
The Baum-Connes conjecture with coefficients is : for all $G$-algebras $A$, $\mu_{G,A}$ is an isomorphism of abelian groups.
\end{conj}

The conjecture was first stated in \cite{BaumConnes} and in \cite{BaumConnesHigson}. The statement is a generalization of the Connes-Kasparov conjecture, which is simply, when one knows the Baum-Connes conjecture, the latter for almost connected groups (i.e. locally compact groups $G$ such that $G/G_0$ is compact, where $G_0$ is the connected component of the identity).\\

Here is the status of the conjecture :
\begin{itemize}
\item[$\bullet$] The Connes-Kasparov conjecture was established by J. Chabert, S. Echterhoff and R. Nest, who proved that $\mu_{G}$ is an isomorphism for every secound countable almost connected group $G$, and every group of $k$-rational points of a linear algebraic group over a local field of characteristic $0$.
\item[$\bullet$] N. Higson and G. Kasparov proved that the conjecture with coefficients holds for groups having Haagerup property. 
\item[$\bullet$] V. Lafforgue proved the conjecture with coefficients for hyperbolic groups.
\item[$\bullet$] V. Lafforgue showed the conjecture for every semi-simple Lie group, real or $p$-adic reductive, and discrete cocompact subgroup of real rank $1$ Lie groups or of $SL(3,k)$ for any local field $k$. \\
\end{itemize}

Recall that a function $h : G\rightarrow \R$ is said to be of negative type if :
\begin{itemize}
\item[$\bullet$] $h_{|G^{(0)}}=0$,
\item[$\bullet$] $h(g)=h(g^{-1})$ for all $g\in G$,
\item[$\bullet$] for all $t_i\geq 0, i=1,..,n$ such that $\sum t_i=0$ and all $g_1,...,g_n\in G$ having the same range, $\sum_i t_i t_j h(g_i^{-1} g_j )\leq 0$.
\end{itemize}
The function $h$ is said to be locally proper if, for every compact $K \subset G^{(0)}$, $h_{|G_K^K}$ is proper.\\

\begin{definition}
A topological groupoid $G$ is said to be a-T-menable if it satisfies one of the following equivalent property :
\begin{itemize}
\item[$\bullet$] there exists a continuous field of Hilbert spaces over $G^{(0)}$ with a proper affine action of $G$,
\item[$\bullet$] there exists a locally proper negative type function on $G$.
\end{itemize} 
\end{definition}

For nice topological groupoids, the conjecture is known to be false in full generality \cite{HigsonLaffSk}, but the following result was shown by J-L. Tu in \cite{TuThese}
\begin{thm}\label{Tu}
Let $G$ be a locally compact $\sigma$-compact Hausdorff groupoid. Then if $G$ is a-T-menable, the Baum-Connes conjecture with coefficients holds for $G$.
\end{thm}

\subsection{The Coarse Assembly map}

The construction of the Coarse Assembly map is very similar to the Assembly map for groupoids, but relies on different functors.\\

In \cite{SkTuYu} is defined the Roe algebra $C^*(X,B)$ of a discrete metric space with bounded geometry $X$ with coefficients in an arbitrary $C^*$-algebra $B$. For $R>0$, define $C_R[X,B]$ to be the involutive subspace of operators $T\in \mathcal L_B(l^2(X)\otimes H\otimes B)$ such that $T_{xy}\in \mathfrak K(H\otimes B)$ and $T_{xy}=0$ as soon as $d(x,y)>R$, and $\C[X,B]= \cup_{R>0} C_R[X,B]$. 

\begin{definition}
The Roe algebra $C^*(X,B)$ is defined as the Hibert $B$-module obtained after the completion of $\C[X,B]$ with respect to the operator norm in $\mathcal L_B(l^2(X)\otimes H\otimes B)$.
\end{definition} 

This construction is functorial : any $*$-homomorphism $\phi: A\rightarrow B$ gives rise to a $*$-homomorphism $\phi_X : C^*(X,A)\rightarrow C^*(X,B)$. Moreover this functoriality extends to $KK$-theory : there exists a natural transformation 
\[\sigma_X : KK(A,B)\rightarrow KK(C^*(X,A),C^*(X,B))\]
which respects the Kasparov product, i.e. $\sigma(z\otimes_D z')=\sigma_X(z)\otimes_{C^*(X,D)} \sigma_X(z')\quad,\forall z\in KK(A,D),z'\in KK(D,B)$.\\

Let $X$ be a coarse space with bounded geometry. The Rips complex is the inductive system of finite dimensional simplicial complexes
\[P_E(X)=\{\eta\in \text{Prob}(X)\text{ s.t. supp }\eta\subset E\}\quad,\forall E\in \mathcal E_X,\]
endowed with the $*$-weak topology. It is an inductive system with respect to inclusion of entourages.\\
 
The coarse homology group of $X$ with coefficients in an arbitrary $C^*$-algebra $B$ is defined as 
\[KX_*(X,B)=\varinjlim_{E\in \mathcal E_X} RK(P_E(X),B)\]
where the limit is taken along the inductive system $\{P_E(X)\}_{E\in\mathcal E_X}$. \\

Let $E\in \mathcal E_X$. Any simplex $\eta\in P_E(X)$ can be written as a finite sum $\eta=\sum_x \lambda_x(\eta)\delta_x$ by boundedness of the geometry. Here, $\delta_x$ is the Dirac measure at $x\in X$. The functions $\lambda_x :\eta\mapsto \lambda_x(\eta)$ are continuous and satifisfy $\sum\lambda_x(\eta) = 1,\forall\eta\in P_E(X)$, so that
\[V_0\left\{\begin{array}{rcl} 
 C_0(P_E(X)) & \rightarrow 	& l^2(X)\otimes C_0(P_E(X)) 		\\ 
 f           & \mapsto 		& (\lambda_x^{\frac{1}{2}}f)_{x\in X} 
\end{array}\right.\] 
is an isometry of $C_0(P_E(X))$-modules, and $V_0V_0^*\otimes id_H$ defines a projection in $l^2(X)\otimes H\otimes C_0(P_E(X))$ with finite propagation, so class $[\mathcal L_E]\in K_0(C^*(X,C_0(P_E(X)))$.\\

\begin{definition}
The Coarse Assembly map $\mu:KX_*(X,B)\rightarrow K_*(C^*(X,B))$ for $X$ with coefficients in $B$ is defined as 
\[\forall z\in RK(P_E(X), B), \mu_{X,B}(z)=[\mathcal L_E]\otimes_{C^*(X,C_0(P_E(X)))} \sigma_X(z).\]
\end{definition}

The Coarse Baum-Connes conjecture is the following claim.\\

\begin{conj}[Coarse Baum-Connes conjecture]
For any coarse space $X$ with bounded geometry, the Coarse Assembly map $\mu_{X,\C}$ is an isomorphism.\\
\end{conj}

Here are some remarks about the status of the conjecture :
\begin{itemize}
\item[$\bullet$] the conjecture is known to hold for any coarse space that admits a coarse embedding into a  separable Hilbert space\cite{Yu2}.
\item[$\bullet$] couterexamples have been constructed \cite{HigsonLaffSk},
\item[$\bullet$] Let $\Gamma$ be a finitely generated group, and $|\Gamma|$ the coarse space associated to the word-length metric (any two left invariant metric on $\Gamma$ are coarsely equivalent). Then the Coarse Baum-Connes conjecture for $|\Gamma|$ implies the Novikov conjecture for higher signatures for $\Gamma$.\\
\end{itemize}

Here is another description of the Coarse Assembly map which will be of some interest for us.\\

Let $E\in \mathcal E_X$ be an entourage. Take a cycle $(H,\pi, T)\in E(C_0(P_E(X)), \C)$, and equip the finite dimensional simplicial complex $P_E(X)$ with a metric which restricts to the spherical metric on simplices, so that $H$ is a $P_E(X)$-module that we can suppose standard non-degenrate. Then 
\[T'=\sum \lambda_x^{\frac{1}{2}} T \lambda_x^{\frac{1}{2}} \] 
is a compact perturbation of $T$ in $\mathcal L(H)$ with bounded propagation, and is invertible modulo $C^*(P_E(X),H)$, so defines a class $[T']\in K_0(C^*(P_E(X),H))$. But $P_E(X)$ and $X$ are coarsely equivalent, the barycentric map $P_E(X)\rightarrow X$ is a coarse equivalence which is covered by $V_0$, so that 
\[(Ad_{V_0})_* : K_*(C^*(P_E(X),H))\rightarrow K_*(C^*(X,l^2(X)\otimes H))\]
is an isomorphism, which does only depend on the coarse class of the barycentric map.\\

With that in mind, one can show that $\mu_{X}([H,\pi,T])=(Ad_{V_0})_*[T']$.\\

The point of the article \cite{SkTuYu} is to show that the Coarse Assembly map is equivalent to the Assembly map for the coarse groupoid $G=G(X)$ with coefficients in the $G$-algebra $l^\infty=l^\infty(X,\mathfrak K)$. \\

More precisely, $Z=P_{\overline E}(G)$ is a proper cocompact $G$-space with ismomorphic fibers $Z_x\simeq P_E(X)$, and $l^\infty_x \simeq \mathfrak K$, so that the inclusion of groupoid $\{x\}\hookrightarrow G $ induces a morphism $\iota_Z : RK^G(Z,l^\infty)\rightarrow KK(Z_x,\mathfrak K)$ which is actually an isomorphism. Recall that one can construct an $*$-isomorphism $\Psi_X: l^\infty \rtimes_r G\rightarrow C^*(X)$, and $\mu_X^{(Z_x)}\circ \iota_Z = (\Psi_X)_*\circ \mu_{G,l^\infty}^{(Z)}$ holds and respects inductive limits over $Z$ to give the following commutative diagram :
\[\begin{tikzcd}
K_*^{top}(\underline E G, l^\infty) \arrow{r}{\mu_{G,l^\infty}}\arrow{d}{\varinjlim \iota_Z} & K_*(l^\infty\rtimes_r G)\arrow{d}{(\Psi_X)_*} \\
KX_*(X)\arrow{r}{\mu_X} & K_*(C^*(X))
\end{tikzcd}\]  
with vertical arrows being isomorphism.\\

A key point is to use how analytical properties of $G$ translate coarse properties of $X$. In particular, 
\begin{itemize}
\item[$\bullet$] in \cite{SkTuYu} is proved that $X$ admits a coarse embedding into Hilbert space iff $G(X)$ is a-T-menable, 
\item[$\bullet$] in \cite{SkTuYu} is proved that $X$ has property A iff $G(X)$ is amenable, 
\item[$\bullet$] in \cite{FinnSellFibred}, M. Finn-Sell shows that if $X$ admits a coarse fibered embedding, then $G_\partial = G(X)_{|\partial \beta X}$ is a-T-menable. 
\end{itemize}
All of these results, combined with theorem \ref{Tu}, give examples of coarse spaces satisfying the Coarse Baum-Connes conjecture and the Boundary Coarse Baum-Connes conjecture.


















\section{Groupoïdes}
\begin{frame}
  \tableofcontents[currentsection]
\end{frame}

\section{Preliminaries on topological groupoids}

\begin{definition}
A topological groupoid is the data of 
\begin{itemize}
\item[$\bullet$] two topological spaces $G$ and $G^{(0)}$, called respectively the space of arrows, and the unit space,
\item[$\bullet$] a topological embedding $e: G^{(0)}\rightarrow G; x\mapsto e_x$ called the unit map,
\item[$\bullet$] source and range maps $r,s : G\rightarrow G^{(0)}$ which are continuous, surjective, and satisfy $r\circ e = s\circ e = Id_{G^{(0)}}$,
\item[$\bullet$] a continuous multiplication map $m : G^{(2)}=G\times_{s,r}G\rightarrow G; (g,g')\mapsto gg'$ such that $(gg')g''=g(g'g'')$ for all triple $g,g',g''$ that are composable, and $g e_{s(g)} = e_{r(g)} g =g$,
\item[$\bullet$] an involutive map $\text{inv} : G\rightarrow G: g\mapsto g^{-1}$ which is continous and such that $s\circ \text{inv} = r$, $r\circ \text{inv} = s$ and $\forall g\in G, gg^{-1}=e_{s(g)} \text{ and }g^{-1}g=e_{r(g)}$. This map is called the inverse map.
\end{itemize} 
\end{definition}

\begin{definition}
Let $G$ be a locally compact groupoid, i.e. a topological groupoid locally compact as a topological space.\\
$G$ is said to be étale if the range map $r:G\rightarrow G^{(0)}$ is a local homeomorphism.\\
$G$ is said to be principal if $r\times s : G\rightarrow G^{(0)}\times G^{(0)}$ is injective, transitive if $r\times s $ is surjective, and proper if $r\times s $ is a proper map. 
\end{definition}

If $G$ is étale, then any $g\in G$ has an open neighborhood $U$ such that $r$ and $s$ are injective on $U$. The element $g$ defines in this way a partial homeomorphism $\alpha_U = r\circ (s_{|U})^{-1} : s(U)\rightarrow r(U)$. Such open sets $U$ where $r$ and $s$ are injective are called bisections. The set of bisections is naturally an inverse semi-group. \cite{paterson}\\

For $U,V\subset G^{(0)}$, we will us the notations $G_U=s^{-1}(U)$, $G^V=r^{-1}(V)$ and $G_U^V=s^{-1}(U)\cap r^{-1}(V)$. If one of the subsets is reduced to a single point, we will use $G_x=G_{\{x\}}$, etc. We will write $G_{|U}$ for $G_U^U$. It is a subgroupoid of $G$.\\

A Haar system on $G$ is a family of borelian measures $(\lambda^x)_{x\in G^{(0)}}$ such that 
\begin{itemize}
\item[$\bullet$] the support of $\lambda^x$ is $G^x$, 
\item[$\bullet$] the map $G^{(0)} \rightarrow \C; x\mapsto \int_{G^x} f d\lambda^x$ is continuous for all $f\in C_c(G)$,
\item[$\bullet$] for all $f\in C_c(G)$ and $g\in G$, $\int_{G^{r(g)}} f(h) d\lambda^{r(g)}(dh)= \int_{G^{s(g)}} f(gh) \lambda^{s(g)}(dh)$ holds.
\end{itemize}

If the groupoid is étale, a natural Haar system is the counting measure on the discrete fibers, i.e. $\lambda^x(A) = |A\cap G^x|$ for all $A\subseteq G$.\\

Here are classical examples of groupoids.

\begin{Expl}[Groups and spaces] Any group $G$ is a groupoid over a point $G^{(0)}=*$. $G$ is étale as a groupoid if it is a discrete group. A topological space is also a groupoid over itself, the only arrows being the units over the points : $G^{(0)} = X = G$, $s=r=id_X$, it is always étale.
\end{Expl}

\begin{Expl}[Pair groupoid] Let $X$ be a set, and consider $G=X\times X$ with base space $G^{(0)}=X$ and projection maps as target and source $s=pr_1, r=pr_2$. Composition is given by $(x,y)(y,z)= (x,z)$, and one can easily check that $e_x=(x,x)$ and $(x,y)^{-1} = (y,x)$. $X\times X$ is étale if $X$ is discrete.
\end{Expl}

\begin{Expl}[Equivalence relations] More generally, consider an equivalence relation $R \subset X\times X$ on a topological space $X$. It naturally has a structure of topological groupoid, as a subgroupoid of $X\times X$. 
\end{Expl}

\begin{Expl}[Pull-back]
Let $X$ be a topological space and $f : X\rightarrow G^{(0)} $ a continuous map. The pull-back groupoid of $G$ along $f$, denoted $G[X]$, is defined as $\{(x,g,y)\in X\times G \times X \text{ s.t. }g\in G_{f(x)}^{f(y)}\}$, with base space $X$. A particular example is given by considering an open cover $\mathcal U=(U_j)$ of $G^{(0)}$. Let $p : \coprod_{j} U_j \rightarrow G^{(0)}$ be the natural projection. Then $G[\mathcal U]$ denotes the pull-back of $G$ along $p$.
%We define $G[\mathcal U]$ as the topological groupoid obtained as the set of triples $(U,g,V)$ such that $s(g)\in U$ and $r(g)\in V$ with base space the groupoid of the equivalence relation defined by $g\sim g'$ if and only if $p(g)=p(g')$. It is convient to represent this groupoid as the set of triples $(U,g,V)$ such that $s(g)\in U$ and $r(g)\in V$
\end{Expl}

\begin{Expl}[Action groupoid] Let $G$ be a group acting on a topological space $X$ by homeomorphisms. Define the groupoid $X\rtimes G$ as the topological space $X\times G$ over the base space $X$, endowed with the maps $s(x,g)=x$ and $r(x,g)=g.x$. The unit map will be given by $x\mapsto (x,e_G)$ and $(x,g)^{-1} = (g.x , g^{-1})$. This groupoid is called the action groupoid. It is étale when $G$ is a discrete group.  
\end{Expl}

These examples show how the notion of topological groupoid interpolates between a group and a space. However, other examples of interesting groupoids abound in analysis and geometry. For example, we will define later a so called Coarse Groupoid associated to any Coarse Space, after \cite{SkTuYu}, whose importance cannot be emphasized enough to our approach to the Coarse Baum-Connes conjecture.\\

Other examples includes the holonomy groupoid of a foliation \cite{connesfoliations}, Lie groupoids \cite{paterson}, etc. One of the reasons groupoids are such famous objects in noncommutative geometry is that we can still do analysis on it. According to the general philosophy of Connes \cite{Connes94}, groupoids encode geometric situations and remain available to analysis, even when the situation is so pathological that classical tools become trivial. For example, when confronted to a group action, considering the convolution algebra of the principal action groupoid can be a substitute to the studying the quotient space. It turns out that, when the action is free and proper, the envelopping $C^*$-algebra is Morita equivalent to $C_0(X/G)$, which represent the quotient space by Gelfand duality. However, when the quotient is %badly behaved
singular, the envelopping $C^*$-algebra still carries usefull information, even if the quotient  is such that $C_0(X/G) = \C$.

%%%%%%%%%%%%%%%%%%%%%%%%%%%%%%%%%%
\subsection{Actions of groupoids}
%%%%%%%%%%%%%%%%%%%%%%%%%%%%%%%%%%

\begin{definition}
A left action of $G$ on a topological space $Z$ is given by a continuous map $p : Z \rightarrow G^{(0)}$, called the anchor map, and a map $\alpha : G\times_{s,p} Z \rightarrow Z $ such that :
\begin{itemize}
\item[$\bullet$] $\alpha(g',\alpha(g,z)) = \alpha(g'g,z)$ whenever $(g,g')\in G^{(2)}$ and $p(z)=s(g)$,
\item[$\bullet$] $p(\alpha(g,z))= r(g)$
\item[$\bullet$] $\alpha(e_{p(z)},z)=z$
\end{itemize} 
We will use the notation $\alpha(g,z) = g.z$ when the action is clear from the context, and $Z$ will then be referred as a (left) $G$-space.\\

An right action of $G$ on a topological space $Z$ is given by a continuous map $p : Z \rightarrow G^{(0)}$, again called the anchor map, and a map $\alpha : Z\times_{p,r} G \rightarrow Z $ such that :
\begin{itemize}
\item[$\bullet$] $\alpha(\alpha(z,g),g') = \alpha(z, gg')$ whenever $(g,g')\in G^{(2)}$ and $p(z)=r(g)$,
\item[$\bullet$] $p(\alpha(z,g))= s(g)$
\item[$\bullet$] $\alpha(z,e_{p(z)})=z$
\end{itemize} We will use the notation $\alpha(g,z) = z.g$ when the action is clear from the context.\\ 
\end{definition}

A left action $(Z,p,\alpha)$ is said to be 
\begin{itemize}
\item[$\bullet$] proper if $id_Z \times \alpha : (g,z) \mapsto (z,g.z) $ is proper as a continous map,
\item[$\bullet$] free if $\alpha(g,z)=z \Rightarrow g=e_{p(z)} $.
\end{itemize}
Similarly, we define these notions for right actions.\\

Let $Z$ be a right $G$-space. Define :
\begin{itemize}
\item[$\bullet$] $(Z\rtimes G)^{(0)} = Z$, $Z\rtimes G = Z\times_{r,p} G$ as a topological space,
\item[$\bullet$] $u_{z}= (z,e_{p(z)})$, $ s(z,g) = z$ and $r(z,g)=z.g$, 
\item[$\bullet$] if $y = x.g$, $(x,g)(y,g')= (x,gg') $ and $(x,g)^{-1} = (x.g,g^{-1})$.
\end{itemize}

These maps define a structure of toplogical groupoid on $Z\rtimes G$. It is called the crossed product groupoid of $Z$ by $G$. It is étale if $G$ is.\\

We present now an important class of $G$-spaces called $G$-simplicial complexes. The reader is referred to \cite{TuBC2} for details. 

\begin{definition}
A map between two topological spaces $f : X\rightarrow Y$ is said to be locally injective if there exists an open cover $\mathcal U$ of $X$ such that, for all $U\in \mathcal U$, $f_{|U}$ is injective.
\end{definition}

\begin{definition} \label{Gcomplex}
Let $n\in\N$. A $G$-simplicial complex of dimension $\leq n$ is a pair $(X,\Delta)$ where :
\begin{itemize}
\item[$\bullet$] $X$ is a locally compact $G$-space, called the space of vertices, such that the anchor map $p : X\rightarrow G^{(0)}$ is locally injective;
\item[$\bullet$] $\Delta$ is a closed $G$-invariant subset of the space of measures on $X$, denoted $M_X$, endowed with the weak $*$-topology. Moreover, $\Delta$ contains only probability measures and satisfies :
\begin{itemize}
\item[$\bullet$] for all $\eta\in\Delta$, there exists $x\in G^ {(0)}$ such that $\text{supp }\eta \subseteq p^{-1}(x)$ and $|\text{supp }\eta|\leq n+1$,
\item[$\bullet$] if $\eta' \in \Delta$ and $\eta\in M_X$ such that $\text{supp }\eta \subseteq \text{supp }\eta'$, then $\eta\in \Delta$.
\end{itemize}
For $\eta\in \Delta$, $\text{supp }\eta$ is called a simplex, or a $j$-simplex when $|\text{supp }\eta | = j$.
\end{itemize}
The complex is said to be typed if there exists a discrete $G$-space $T$ and a $G$-equivariant continuous map $\tau : X\rightarrow T$ such that, for every simplex $S$, $\tau_{|S}$ is injective.  
\end{definition}

For any typed proper $G$-compact $G$-simplicial complex $(X,\Delta)$ of dimension $\leq n$, one can decompose it into its $n$-skeleton $Z_0\subseteq Z_1 \subseteq ... \subseteq Z_n$, where each $Z_j$ is a closed $G$-invariant subset of $\Delta$ such that :
\begin{itemize}
\item[$\bullet$] for all $\eta \in Z_j \setminus Z_{j-1}$, $|\text{supp }\eta|= j$,
\item[$\bullet$] $Z_j \setminus Z_{j-1}$ is $G$-equivariantly homeomorphic to $\mathring\sigma_j \times \Sigma_j$, where $\mathring\sigma_j $ is the interior of the standard simplex of dimension $j$, and $\Sigma_j$ is the discrete subspace of centers of $j$-simplices.
\end{itemize}

\begin{definition}
Let $(X,\Delta)$ be a $G$-simplicial complex of dimension $\leq n$. Its barycentric subdivision is the $G$-simplicial complex $(S,\Delta_S)$ of dimension $\leq n$ defined by :
\[S = \cup_{\eta\in \Delta}\text{supp } \eta\subseteq \Delta ,\]
i.e. the space of vertices $S$ contains $X$ plus the barycentric center of any simplex of $\Delta$, and $\Delta_S $ satisfies that $\sigma = \{ \eta_0, ...,\eta_k\}$ is a simplex in $\Delta_S$ iff $\cup_{0 \leq j \leq k}\text{ supp }(\eta_j)$ is a simplex of $\Delta$.
\end{definition}

This construction shows that, up to replacing a $G$-simplicial complex of dimension $\leq n$ by its barycentric subdivision, we can always suppose that it is a typed $G$-simplicial complex of dimension $\leq n$.

%%%%%%%%%%%%%%%%%%%%%%%%%%%%%%%%%%%%%%%%%%%%%%%
\subsection{$G$-algebras and crossed products}
%%%%%%%%%%%%%%%%%%%%%%%%%%%%%%%%%%%%%%%%%%%%%%%

This paragraph describes what is an action of a topological groupoid on a $C^*$-algebra, and how to construct the reduced crossed-product out of the action.

\begin{definition}
Let $X$ be a topological space. A $C(X)$-algebra is a $C^*$-algebra together with a $*$-homomorphism $\theta : C_0(X)\rightarrow Z(\mathcal M(A))$ such that $\overline{\theta (C_0(X)) A} = A$. The couple $(A,\theta)$ is then called a $C(X)$-structure.
\end{definition} 

If $f : X\rightarrow Y$ is a continuous map between locally compact spaces, and $A$ a $C(Y)$-algebra, then $f$ induces a $C(X)$-structure on $A$ by $f^* : C(Y) \rightarrow C_b(X); g\mapsto g\circ f$. Define $f^* A = A \otimes_{f^*} C_0(X)$ the resulting $C(X)$-algebra. For details on tensor products of $C(X)$-algebras, see \cite{LeGall} for instance. As a particular case, one gets the notion of a fiber of a $C(X)$-algebra : for $x\in X$, $A_x = ev_x^* \ A$ is called the fiber over $x$, where $ev_x$ is the inclusion of $x$ in $X$. It is convenient to see $A_x$ as $A/ (I_x .A)$ where $I_x$ is the ideal ker $(ev_x : C_0(X) \rightarrow \C)$. Moreover we can describe the fibers of $f^*A$ : $(f^*A)_x \cong A_{f(x)}$ for all $x\in X$.\\

A homomorphism between two $C(X)$-algebras $\alpha : A\rightarrow A'$ is a $*$-homomorphism which commutes with the action of $C_0(X)$, i.e. $\alpha(\theta(f)a)=\theta'(f)\alpha(a)$ for all $a\in A,f\in C_0(X)$. Its fibers are defined unambiguously by $\alpha_x = id_\C \otimes_{ev_x} \alpha$.\\

Notice that, if $A$ is a $C(G^{(0)})$-algebra, then $s^* A = A \otimes_s C_0(G)$ and $r^* A = A \otimes_r C_0(G)$ are $C(G)$-algebra with fibers $(s^*A)_g \cong A_{s(g)}$ and $(r^*A)_g \cong A_{r(g)}$ respectively, for all $g\in G$.

\begin{definition}
An action of a groupoid $G$ on a $C^*$-algebra $A$ is a triple $(A,\theta,\alpha)$ where $(A,\theta)$ is a $C(G^{(0)})$-structure and $\alpha : s^* A \rightarrow r^* A$ is an isomorphism of $C(G)$-algebras such that $\alpha_{gg'}= \alpha_g\circ \alpha_{g'}$ for all $(g,g')\in G^{(2)}$.
\end{definition} 

If $A$ is a $G$-algebra, and $U \subset G^{(0)}$ is an open subset, recall from \cite{LeGall} that the restriction $A_U$ is defined as $\eta^* A = C_0(U) \otimes_\eta A$ where $\eta : C(U) \hookrightarrow C_0(G^{(0)})$ is the canonical inclusion. 
%Let $F= G^{(0)}\backslash U$. As $A_U$ is naturally embedded into $A$, define $A_F$ to be the complementary of $A_U$. \\

The space of continuous sections with compact support is defined as :
\[C_c(G,A) = \cup_{U\in \mathcal E} A_U,\]
where $U$ runs through the open relatively compact subsets of $G$. It is naturally an involutive algebra with the convolution product
\[\phi\ast \psi (g) = \int_{G^{r(g)}} \phi(h)\psi(h^{-1}g)\lambda^{r(g)}(h).\] 

\begin{rk}
Another possible definition of $C_c(A,G)$ would be to take a tensor product of Fréchet algebras $C_c(G)\otimes A$, which gives a Fréchet algebra. Even if this definition seems more elegant, this would require more definitions on Fréchet $G$-algebras, which we do not want to dive into.  
\end{rk}

%Let $E$ be a compact subset of $G$. The space of continuous sections with support in $E$ is defined as $A_E$. Then 
%\[C_c(G,A) = \cup_{E\in \mathcal E} A_E\]
%is called the space of continuous sections with compact support. It is naturally a $*$-algebra with the convolution product
%\[\phi\ast \psi (g) = \int_{G^{r(g)}} \phi(h)\psi(h^{-1}g)\lambda^{r(g)}(h).\] 

%If $f\in C_c(G,A)$, define $\lambda(f)\in \mathcal L(L^2(G,A))$ as the convolution operator 
%\[\lambda(f)\eta = f\ast \eta\]
%and the reduced norm as $||f||_r= ||\lambda(f)||_{\mathcal L(L^2(G,A))}$.

Let $g\in G$ and $f\in C_c(G,A)$. Choose an open relatively compact subset $U\subseteq G$ which contains supp $f$. We can define $f(g)$ as $1\otimes_{ev_g} f \in \C \otimes_{ev_g} C_0(U)\otimes_\eta A\cong A_{s(g)}$. This jusitifies the name "continuous sections" for $C_c(G,A)$. \\

Let $A$ be a $G$-algebra and $x\in G^{(0)}$. As the groupoid $G$ is étale, $G_x$ is discrete and $l^2(G_x)$ is a separable Hibert space with dimension the cardinality of $G_x$. Define $l^2(G_x,A)$ as the $A_x$-Hilbert module $l^2(G_x)\otimes A_x$. We denote by $||\cdot||_2$ the norm induced by the inner product on $l^2(G_x,A)$ i.e. 
\[ \begin{array}{c}
\langle \xi,\eta\rangle_{A_x}= \sum_{g\in G_x} \xi(g)^* \eta(g),\\
||\xi||_2= ||\langle \xi,\xi \rangle||^\frac{1}{2} 
\end{array}\] 
for all $\xi,\eta\in l^2(G_x,A)$.\\

Define, for $f\in C_c(G,A)$, the reduced norm as 
\[||f||_r= \sup_{x\in G^{(0)}}\{ ||f\ast \xi||_2 : \xi\in l^2(G_x,A) \text{ s.t. } ||\xi||_2\leq 1\}.\]


\begin{definition}
The reduced crossed product $A\rtimes_r G$ is defined as the completion of $C_c(G,A)$ under the reduced norm.
\end{definition}

% Propriete, fonctorialite etc....

\subsection{Equivariant Hilbert modules}
% A FINIR
Let $B$ be a $G$-algebra, with action $\beta$, and $E$ a $B$-Hilbert module. As earlier, we can define the fiber $E_x$ of $E$ over $x\in G^{(0)}$. Define $s^* E = E \otimes_s C_0(G)$ and $r^* E = E \otimes_r C_0(G)$ which are $s^*B$ and $r^*B$-Hilbert modules respectively, with fibers $(s^* E)_g \cong E_{s(g)}$ and $(r^* E)_g \cong E_{r(g)}$ for all $g\in G$. We can endow $r^*E$ with a structure of $s^* B$-Hilbert module with the help of the isomorphism $\beta^{-1} : r^* B\rightarrow s^* B$.

\begin{definition}
An equivariant Hilbert $B$-module is a Hilbert $B$-module $E$ together with an unitary $V\in\mathcal L_{s^* B}(s^*E ,r^* E)$ such that $V_gV_{g'} = V_{gg'}$ for all $(g,g')\in G^{(2)}$. If the underlying $G$-algebra is clear, we will often abbreviate $G$-equivariant Hilbert $B$-module by $G$-module.
\end{definition}

\begin{Expl} Define $L^2(G)$ as the completion of $C_c(G)$ seen as a pre-Hilbert $C_0(G^{(0)})$-module with the $C_0(G^{(0)})$-valued inner product 
\[\langle \xi ,\eta\rangle (x) = \sum_{g\in G^x} \xi(g)^*\eta(g) \quad \forall \xi,\eta\in C_c(G).\]
Its fibers are isomorphic to $l^2(G^x)$, and the action $\lambda : s^*L^2(G) \rightarrow r^* L^2(G)$ is given by left translation, i.e. if $\xi\in l^2(G^x)$, then
\[(\lambda_g\xi)(h) = \xi(g^{-1}h)\quad \forall g,h\in G^x.\]
\end{Expl}

\begin{Expl} Let $B$ be a $G$-algebra, the $C(G^{(0)})$-structure being given by $\theta : C_0(G^{(0)})\rightarrow Z(\mathcal M(B))\subseteq \mathcal L_B(B)$. The $G$-module $L^2(G,B)$ is defined as the tensor product $ L^2(G)\otimes_\theta B$. Its fibers are $l^2(G^x,B)$ and the action is given by $\lambda\otimes id_{B}$.
\end{Expl}

\begin{Expl} More generally, let $B$ be a $G$-algebra and $E$ be a $B$-module, with $C(G^{(0)})$-structure given by $\theta : C_0(G^{(0)})\rightarrow Z(\mathcal M(B))\subseteq \mathcal L_B(B)$. The $G$-equivariant $B$-Hilbert module $L^2(G,E)$ is defined as the tensor product $ L^2(G)\otimes_\theta E$.
\end{Expl}

\begin{rk}
We could alternatively have defined, for $A$ a $G$-algebra, $A\rtimes_r G$ as a subalgebra of $\mathcal L_A(L^2(G,A))$. Indeed, any $f\in C_c(G,A)$ acts as a convolution operator on $L^2(G,A)$ :
\[(\lambda(f)\xi)(h) = \sum_{g\in G^x} f(g)\alpha_g(\xi(g^{-1} h))\quad \forall \xi\in L^2(G,A),g,h\in G^x,\]
and $||f||_{r} = ||\lambda(f)||_{\mathcal L_A (L^2(G,A))}$.
\end{rk}

\begin{lem}[Equivariant stabilization lemma \cite{LeGall}] \label{GStabilization} Let $E$ be a countably generated $G$-equivariant $B$-Hilbert module, then there exists an $G$-equivariant isomorphism of Hilbert module $E \oplus H_B \cong H_B$.
\end{lem}

%%%%%%%%%%%%%%%%%%%%%%%%%%%%%%%%%%%%%%
\subsection{Equivariant $KK$-theory}
%%%%%%%%%%%%%%%%%%%%%%%%%%%%%%%%%%%%%%

From now on, all the groupoids considered are locally compact and $\sigma$-compact.\\

In \cite{LeGall}, Le Gall defined, for any locally compact groupoid $G$, and $G$-algebras $A$ and $B$, $\Z_2$-graded abelian groups $KK^G(A,B)$ generalizing Kasparov's equivariant $KK$-theory for a locally compact group. Elements of $KK^G(A,B)$ are defined as equivalence classes of triples $(E,\pi,T)$ that we will call $K$-cycles.

\begin{definition} A $K$-cycle is a triple $(E,\pi,T)$ where :
\begin{itemize}
\item[$\bullet$] $E$ is a right $\Z_2$-graded $B$-Hilbert module, equipped with a unitary action of $G$, i.e. an even unitary $V\in\mathcal L_B(s^*E,r^*E)$ such that $V_g V_{g'} = V_{gg'}$ for all $(g,g')\in G^{(2)}$,
\item[$\bullet$] $\pi : A\rightarrow \mathcal L_B(E)$ is a $*$-homomorphism whose image is contained in even operators, and we will often write $a$ instead of $\pi(a)$,
\item[$\bullet$] $T\in\mathcal L_B(E)$ is an odd bounded $B$-linear operator satifying the $K$-cycle relations i.e. $[a, T]$, $a(T-T^*)$, $a(T^2-1)$ are compact operators in $\mathfrak K_B(E)$ and $a(r^* T -V s^*T V^*)$ is a compact operator in $\mathfrak K_{r^* B}(r^* E)$ for all $a\in A$.
\end{itemize} 
The set of such $K$-cycles is denoted by $\mathbb E^G(A,B)$.
\end{definition}

\begin{definition}
Two $K$-cycles $(E,\pi,T)$ and $ (E',\pi',T')$ in $\mathbb E^G(A,B)$ are homotopic if there exists $ (\tilde E,\tilde \pi,\tilde T) \in\mathbb E^G(A,B[0,1])$ such that $(\tilde E_0,\tilde \pi_0,\tilde T_0)\simeq (E,\pi,T)$ and $(\tilde E_1,\tilde \pi_1,\tilde T_1)\simeq (E',\pi',T')$. The action on $B[0,1]$ is the action obtained as the tensor product of the action on $B$ and of the trivial action on $C[0,1]$.
\end{definition}

\begin{definition} $KK_0^G(A,B)$ is defined as the quotient of $\mathbb E^G(A,B)$ by the equivalence relation given by homotopy. Let $\C_1$ denotes the trivial $G$-Hilbert module $\C^2$ with the standard grading given by $\begin{pmatrix}0 & 1 \\ 1 & 0 \end{pmatrix}$. $KK_1^G(A,B)$ is defined as $KK_0^G(A,B \otimes \C_1)$. 
\end{definition}

Doing computations, we will often use different equivalence relations which turn out to coincide when $A$ is separable and $B$ is $\sigma$-unital. See \cite{blackadar} for details.

\begin{definition} Unitary and operator homotopy equivalences.
\begin{itemize}
\item[$\bullet$]  Unitary equivalence is defined as $(E,\pi,T)\sim_{u} (E',\pi',T')$ if there exists a unitary $u\in\mathcal L_B(E,E')$ such that $T = u^* T' u$ and $\pi(a) = u^* \pi'(a) u , \forall a\in A$.
\item[$\bullet$] Let $\mathbb D^G(A,B) = \{(E,\pi,T) \in\mathbb E^G(A,B) \text{ s.t. } \pi(a)(F^2-1) = 0, [\pi(a),F]=0 \text{ and }V(s^*T)V^* = r^*T,\forall a\in A \}$. It is called the set of degenerate $K$-cycles.   
\item[$\bullet$] An operator homotopy of $K$-cycles is a family $[0,1]\rightarrow \mathbb E^G(A,B) ; s\mapsto (E,\pi,T_s)$, where $s\mapsto T_s$ is norm-continuous. 
\item[$\bullet$] $\sim_{oh}$ is the equivalence relation on $\mathbb E^G(A,B)$ generated by operator homotopy and addition of degenerate elements, i.e. $(E,\pi,T)\sim_{oh} (E',\pi',T')$ if there exists $(E'',\pi'',T'')\in\mathbb E^{G}(A,B)$ such that $(E,\pi,T)\oplus (E'',\pi'',T'')$ and $(E',\pi',T')\oplus (E'',\pi'',T'')$ are operator homotopic, up to unitary equivalence.
\end{itemize}
\end{definition}

The following proposition summarizes propositions $17.2.3$ and $17.2.5$ from \cite{blackadar}. 
 
\begin{prop}[\cite{blackadar}] If $(E,\pi,T)\in\mathbb D^G(A,B)$, then $(E,\pi,T)$ is homotopic to the $0$ $K$-cycle.
If $(E,\pi,T')$ is a compact pertubation of $(E,\pi,T)$ in $\mathbb E^G(A,B)$, i.e. $T'-T\in \mathfrak K_B(E)$, then $(E,\pi,T')$ is operator homotopic to $(E,\pi,T)$.\\
Moreover, if $A$ is separable and $B$ is $\sigma$-unital, then $\sim_{oh}$ coincides with homotopy equivalence of $K$-cycles.
\end{prop}

One of the properties that makes equivariant $KK$-theory so powerful is the existence of a product, called the Kasparov product. Its existence is highly non trivial, and it fails to exists in the non-separable case or the $KK$-theory for Banach algebras developed by V. Lafforgue \cite{Lafforgue} for example.

\begin{prop}\cite{LeGall} Let $A$, $B$ and $C$ be $G$-algebras such that $A$ is separable. There exists a map 
\[\otimes_B :\left\{\begin{array}{ccc} KK^G(A,B)\times KK^G(B,C) & \rightarrow & KK^G(A,C) \\ (z,z') & \mapsto & z\otimes_B z' \end{array}\right. \]
which is associative and distributive w.r.t. the addition of $K$-cycles. The Kasparov product is covariant w.r.t. $C$ and contravariant w.r.t. $A$.
\end{prop}

\begin{rk}
By the equivariant stabilization lemma \ref{GStabilization}, any $B$-Hilbert module $E$ sits as a complemented $G$-submodule of the standard module $H_B$, and as $[E,\pi,T] = [E \oplus H_B,\pi\oplus 0,T\oplus 0]$, we can suppose $E = H_B$.
\end{rk}
% KASPAROV PRODUCT

Here are some examples of $K$-cycles. For detailed proofs of the various constructions, we send the reader to \cite{blackadar} and \cite{LeGall}.

\begin{Expl}
If $\phi : A\rightarrow B$ is a $G$-equivariant $*$-homomorphism, it induces a $K$-cycle $(B,\phi,0)\in \mathbb E^G(A,B)$ whose class is denoted by $[\phi]\in KK_0^G(A,B)$. This induces a covariant functor $\phi_* = - \otimes_A [\phi] : KK^G(D,A)\rightarrow KK^G(D,B)$ and a contravariant functor $\phi^*=  [\phi] \otimes_B - : KK^G(B,D)\rightarrow KK^G(A,D)$ for any $G$-algebra $D$. Moreover, if $\phi' : B\rightarrow C$, the following relation holds 
\[ [\phi'\circ \phi] = [\phi]\otimes_B [\phi'].\]
\end{Expl}

\begin{Expl} Let $0 \rightarrow B\otimes \mathfrak K \rightarrow_\iota E \rightarrow_p A \rightarrow 0 $ be a $G$-equivariant split exact sequence of $G$-algebras. Denote by $s : A \rightarrow E$ the cross section of $p$ and by $\sigma $ the compostion $s\circ p$. Then 
\[(H_B \oplus H_B^{op},\omega \oplus (\omega \circ \sigma ), \begin{pmatrix} 0 & 1 \\ 1 & 0\end{pmatrix} )\] 
defines a $K$-cycle, hence an element of $KK_0^G(A,B)$ called the splitting morphism of the extension. % Definir w et op
\end{Expl}

The first example allows one to think of $KK(A,B)$ as generalized morphisms between $A$ and $B$ : every $*$-homomorphism induces a $K$-cycle such that composition is preserved. Actually, it turns out that generalized morphisms are not so far from coming from genuine $*$-homomorphisms. Indeed, H. Oyono-Oyono has shown in the appendice of \cite{LaffOY} that every element of $KK^G(A,B)$ can be written as the Kasparov product of at most $d$ elements, each one coming either from a $*$-homormorphism or from a $KK$-inverse of a $*$-homomorphism.\\

%%%%%%%%%%%%%%%%%%
%%  PROPERTY D  %%
%%%%%%%%%%%%%%%%%%

More precisely, let us define decomposition property $(d)$.

\begin{definition}
Let $d$ be a positive integer. An element $z\in KK^G(A,B)$ is said to satisfy decomposition property $(d)$ if
\begin{itemize}
\item[$\bullet$] there exist $G$-algebras $A_0$, $A_1$, ..., $A_d$ such that $A_0=A$ and $A_d=B$, 
\item[$\bullet$] there exist elements $z_j \in KK^G(A_{j},A_{j+1})$ for $j\in\{0,..,d-1\}$ such that, either $z_j$ is induced by a $G$-morphism $A_j \rightarrow A_{j+1}$, either there exists a $G$-morphism $\phi_j : A_{j+1}\rightarrow A_j$ such that $z_j \otimes_{A_{j+1}} [\phi_j] = 1_{A_j}$ and $ [\phi_j] \otimes_{A_{j}} z_j  = 1_{A_{j+1}}$,
\end{itemize}
such that $z = z_1 \otimes_{A_1}  ... \otimes_{A_{d-1}} z_{d-1} $ holds.
\end{definition}

Then, the following theorem is true for a universal constant $d$, which does not depend on the groupoid. It will be crucial to prove that the controlled Kasparov and Roe transforms, applications to be defined later, respect the Kasparov product. 

\begin{thm}[\cite{LaffOY}]\label{propertyD}
Let $G$ be a locally compact groupoid with Haar system. Then every element $z\in KK^G(A,B)$ has decomposition property $(d)$.
\end{thm}

We saw that $KK^G(A,B)$ is a bifunctor w.r.t. $*$-homomorphisms, contravariant in the $A$ variable and covariant in the $B$ variable. Let us describe how $KK^G$ transforms under morphisms of groupoids. See $\cite{LeGall}$ for a proof.

\begin{prop} Let $G$ and $G'$ be étale groupoids and $\phi :G\rightarrow G'$ a strict morphism of topological groupoids, i.e. $\phi$ is continuous such that $\phi(e_x)=e_{\phi^{(0)}(x)}$ and $\phi(g)\phi(g') = \phi(gg')$ for all $(g,g')\in G^{(2)}$. Let $A$ and $B$ be $G'$-algebras, and $\phi^* A = A\otimes_\phi C_0(G)$ and $\phi^* B = B\otimes_\phi C_0(G)$ the pull back $G$-algebras. Then $\phi$ induces a homomorphism of $\Z_2$-graded abelian groups
\[\phi^* : KK^{G'}(A,B)\rightarrow KK^G(\phi^* A,\phi^* B),\]
such that, if $\phi' : G'\rightarrow G''$ is another strict morphism, then $(\phi'\circ\phi)^* =\phi^*\circ(\phi')^*$. In other words, $KK^G$ is contravariant in the $G$ variable.
\end{prop}

The next two propositions reminds the reader about crucial natural morphisms built by P.Y. Le Gall in his thesis that we will need later on. We only state results and remind the reader about the constructions, without proofs. The reader can consult \cite{LeGall} for details.

\begin{prop}\cite{LeGall} Let $A$, $B$, and $D$ be $G$-algebras. Then 
\[\tau_D : 
\left\{\begin{array}{rcl} 
KK^G(A,B) & \rightarrow & KK^G(A\otimes D,B\otimes D) \\ 
\ [E,\pi,T] & \mapsto     & [E\otimes_{C_0(G^{(0)})} D,\pi\otimes_{C_0(G^{(0)})} id_D,T\otimes_{C_0(G^{(0)})} id_D]
\end{array}\right. \] 
defines a homomorphism of $\Z_2$-graded abelian groups which respects the Kasparov product, i.e. 
\[\tau_D(x\otimes_B y) = \tau_D(x)\otimes_{B\otimes D} \tau_D(y)\quad \forall x\in KK^G(A,B),y\in KK^G(B,C).\]
\end{prop}

The next transformation is called the descent functor, or the Kasparov transform.

\begin{prop}\cite{LeGall} Let $A$, $B$, $C$ and $D$ be $G$-algebras. Then 
\[j_G : 
\left\{\begin{array}{rcl} 
KK^G(A,B) & \rightarrow & KK^G(A\rtimes_r G,B\rtimes_r G) \\ 
\ [E,\pi,T] & \mapsto     & [E_G,\pi_G,T_G]
\end{array}\right. \] 
defines a homomorphism of $\Z_2$-graded abelian groups which respects the Kasparov product, i.e. 
\[j_G(x\otimes_B y) = j_G(x)\otimes_{B\rtimes_r G} j_G(y)\quad \forall x\in KK^G(A,B),y\in KK^G(B,C).\]
\end{prop}

The remainder of the section is devoted to compile some useful facts that we will use.\\

We will use the following result.\cite{blackadar}\cite{CuSk}

%\begin{prop}
%Let $z\in KK^G_1(A,B)$. The induced homomorphism $-\otimes_A z : K_*(A)\rightarrow K_*(B)$ can be realized, up to Morita equivalence, as the boundary of the following semi-split $G$-equivariant extension of $G$-algebras
%\[\begin{tikzcd}[column sep = small] 
%0 \arrow{r}& \mathfrak K \otimes B \arrow{r} & E^{\pi,T} \arrow{r} & A \arrow{r} & 0\end{tikzcd}\]
%where $E^{\pi,T}=\{(a,P\pi(a)P +y) : a\in A , y\in \mathfrak K \otimes B\}$, the arrows are the obvious inclusion and projection, and the completely positive section is $a\mapsto (a, P \pi(a) P)$. Here, we implicitely choose a $K$-cycle $(H_B,\pi,T)$ representing $z$, and $P=\frac{1+T}{2}$. The fact that this choice does not affect the boundary will be proven in the upcoming sections. 
%\end{prop}

\begin{prop}
Let $z\in KK^G_1(A,B)$ and choose any $K$-cycle $(H_B,\pi,T)$ representing $z$. Let $P=\frac{1+T}{2}$. The induced homomorphism 
\[-\otimes_A z : K_*(A)\rightarrow K_*(B)\]
can be realized, up to Morita equivalence, as the boundary of the following semi-split $G$-equivariant extension of $G$-algebras
\[\begin{tikzcd}[column sep = small] 
0 \arrow{r}& \mathfrak K \otimes B \arrow{r} & E^{(\pi,T)} \arrow{r} & A \arrow{r} & 0\end{tikzcd}\]
where $E^{\pi,T}=\{(a,P\pi(a)P +y) : a\in A , y\in \mathfrak K \otimes B\}$, the arrows are the obvious inclusion and projection, and the completely positive section is $a\mapsto (a, P \pi(a) P)$. The fact that the choice of the representent does not affect the boundary will be proven in the upcoming sections. 
\end{prop}

The following result will be constantly used in the upcoming sections. In short, it explains how to construct the boundary map associated to any extension $0\rightarrow J\rightarrow A \rightarrow A/J \rightarrow 0$ as the Kasparov product of an element $[\partial_{J,A}]\in KK^G_1(A/J,J)$. For references, see \cite{blackadar}.\\

Recall that, for any $G$-equivariant $*$-homomorphism $\phi : A \rightarrow B$, we can define its cone denoted $C_\phi$ by 
\[C_\phi = \{(a,f)\in A\times CB \text{ s.t. } \phi(a)=f(0)\}.\]
If any $G$-algebra $A$ is given, $0\rightarrow SA\rightarrow CA\rightarrow A \rightarrow 0$ is a $G$-equivariant semi-split extension, and defines an element $[\partial_A]=\tau_A([\partial_\C])\in KK^G_1(A,SA)$ which is invertible. Now, give yourself a semi-split extension 
\[\begin{tikzcd}[column sep = small]0 \arrow{r} & J \arrow{r}{\iota} & A \arrow{r}{q} & A/J \arrow{r}& 0 \end{tikzcd}\] 
and define $\alpha : J\rightarrow C_q ; x\mapsto (x,0)$ and the canonical inclusion $\beta : S (A/J)\hookrightarrow C_q$. It turns out that $[\alpha]\in KK_0^G(J,C_q)$ is invertible.

\begin{prop}
The boundary map $\partial_{J,A} : K_*(A/J)\rightarrow K_{*+1}(J)$ is given by the Kasparov product by an element $[\partial_{J,A}]\in KK^G_1(J,A/J)$ which satisfies 
\[[\partial_{J,A}] = [\partial_{A/J}]\otimes_{} \beta^*([\alpha]^{-1}).\]
\end{prop}

The next lemma will be used during the proof of the functorial properties of the controlled Roe and Kasparov tranformations that will be defined later. These transformations are defined on $K$-cycle which are put in standard form, i.e. $[H_B,\pi,T]$. Changing the coefficients leads to $K$-cycles which are not of this form, hence the lemma. It states that we can always describe explicitely the standard form of a $K$-cycle pushed forward.

\begin{lem}\label{isometryKK}
Let $\phi : B\rightarrow B'$ be a $G$-equivariant homomorphism, and $z=[H_B,\pi, T]\in KK^G(A,B)$. Let $V\in\mathcal L_{B'}(H_B\otimes B', H_{B'})$ be the isometry of the remark \ref{isometry} and $p = VV^*\in\mathcal L_{B'}(H_{B'})$. Define $\pi' : A\rightarrow \mathcal L_{B'}(H_{B'})$ as $\pi'(a) = V\pi(a)V^*$ and $T'= V(T\otimes_\phi 1)V^* + 1-p \in \mathcal L_{B'}(H_B')$. Then $(H_{B'},\pi',T')$ is a $K$-cycle and 
\[g^*(z) = [H_B\otimes B',\pi\otimes_\phi 1,\phi_*(T)]=[H_B', \pi', T']\text{ in } KK^G(A,B').\]
\end{lem} 

\begin{dem}
As a simple computation shows : 
\[\begin{array}{ll} [\pi'(a),T'] = V[\pi(a),T]V^*, &  T'^*-T' = V(T^*\otimes_\phi 1-T\otimes_\phi 1)V^* \\
 T'^2-T' = V(T^2\otimes_\phi 1-T\otimes_\phi 1)V^* & \text{and } r^* T' - U s^*T' U^* = V((r^*(T) - U s^*(T) U^*)\otimes_\phi 1 ) V^*
\end{array} \]
(where $U$ is the action of $G$ on $E$). Hence $[H_{B'},\pi',T']$ satisfies the relations of $K$-cycle.\\

Remark \ref{isometry} shows that $ (H_B\otimes_\phi B' )\oplus (1-p) H_{B'} \cong H_{B'}$.\\

Put $E= H_{B'} \oplus (H_B\otimes_\phi B')$ and $R_s =\begin{pmatrix} \cos (s) id_{H_{B'}} & -\sin (s) V \\  \sin (s)V^* & \cos (s) V^*V \end{pmatrix}$ , and define
\[\begin{array}{c} 
T_s = R_s\begin{pmatrix} V(T\otimes_\phi 1)V^* +1-VV^* & 0 \\ 0 & -T\otimes_\phi 1 \end{pmatrix}R_s^*,\\
\pi_s(a)= \begin{pmatrix} V(\pi(a)\otimes_\phi 1) V^* & 0 \\ 0 & \pi(a)\otimes_\phi 1\end{pmatrix}
\end{array}.\]
Then $ s\mapsto [E, \pi_s,T_s ] $ is an homotopy from $[H_{B'}, \pi', T']-[H_B\otimes B',\pi\otimes_\phi 1,\phi_*(T)]$ to the cycle given by 
\[\begin{pmatrix} V(T\otimes_\phi 1)V^*  & 0 \\ 0 & -T\otimes_\phi 1 \end{pmatrix}\]
which is $0\in KK^G(A,B)$. \\
\qed\end{dem}








































\section{Applications d'assemblage}
\begin{frame}
  \tableofcontents[currentsection]
\end{frame}
\section{Assembly Maps}

In this section, all groupoids $G\rightrightarrows G^{(0)}$ are étale. 

\subsection{Proper groupoids and proper actions}

\begin{definition}
A topological groupoid $G\rightrightarrows G^{(0)}$ is proper if the map $G\times_s G^{(0)}\rightarrow G^{(0)}\times G^{(0)}$ defined by $(g,z)\mapsto (z,gz)$ is a proper map.\\
A $G$-space $Z$ is said to be proper if the crossed-product groupoid $Z\rtimes G$ is proper. If the space of orbits $Z/G$ is compact, $Z$ is called $G$-compact.\\
\end{definition}

\begin{definition}
A $G$-space $\underline E G$ is called universal if, given any proper $G$-compact $G$-space $Z$, there exists a continuous $G$-equivariant map $Z\rightarrow \underline E G$.
\end{definition}

We now construct a model for universal spaces of $G$, based on J.-L.Tu's work \cite{TuBC2}. For any compact subset $K\subset G$, define $P_K(G)$ to be the space of probability measures $\nu $ with support contained in one and only one fiber $G^x$ for some $x\in G^{(0)}$, and such that if $g,g'\in \text{supp }(\nu)$, then $g'g^{-1}\in K$. We endow $P_K(G)$ with the weak-$*$ topology.\\

The action of $G$ is defined by translation. The momentum map $P_K(G)\rightarrow G^{(0)}$ is just the map associating to $\nu$ the only $x$ such that $\text{supp }(\nu) \subset G^x$. As the fibers are discrete, any $\nu\in P_K(G)$ can be represented as as sum $\nu = \sum_{g\in G^x} \lambda_g(\nu)\delta_g$, where $\delta_g$ is the Dirac measure at $g\in G $. The continuous functions $\lambda_g$ are called coordinate functions and satisfy $\sum_{g\in G^x} \lambda_g(\nu) =1$ for every $x\in G^{(0)}$. The action of $G$ is given by $g\lambda_h = \lambda_{g^{-1}h}$.\\

\begin{lem}[Tu,\cite{TuBC2}]
The action of $G$ on $P_K(G)$ is proper and cocompact.
\end{lem}  

\begin{lem}[Tu,\cite{TuBC2}]
Let $Z$ be a proper $G$-compact $G$-space. Then there exists a compact subset $K\subset G$ and a $G$-equivariant continuous map $Z\rightarrow P_K(G)$.
\end{lem}  

\subsection{Equivariant $K$-homology}

We will use the equivariant $KK$-theory developped by Le Gall in his thesis \cite{LeGall}, which is an extension of the usual equivariant $KK$-theory of Kasparov. Recall that, if $A$ and $B$ are two $G$-algebras, elements of $KK^G(A,B)$ are homotopy classes of triple $(E,\pi,T)$ where :\\

\begin{itemize}
\item[$\bullet$] $E$ is a $G$-module,
\item[$\bullet$] $\pi : A\rightarrow \mathcal L_B(E)$ is a $*$-homomorphism,
\item[$\bullet$] $T\in \mathcal L_B(E)$ is an adjoinable operator such that the triple satisfies the condition of $K$-cycle : $\pi(a)(T^2-T),\pi(a)(T^*-T)[\pi(a),T]$, and $\pi(a)(T-g.T)$ are compact operators if $E$, for all $a\in A, g\in G$.\\
\end{itemize}

\begin{definition}
Let $Y$ be a proper $G$-space and $B$ a $G$-algebra. Then the analytic $K$-homology of $Y$ with coefficients in $B$ is defined by 
\[RK^G(Y,B) = \varinjlim_{Z\subset Y} KK^G(C_0(Z),B), \]
the inductive limit being taken on proper $G$-compact $G$-subspaces $Z$ of $Y$.
\end{definition}

The previous lemmas assures that this inductive limit can be somehow restricted when it comes to the analytic $>K$-homology of any universal space $\underline E G$ :
\[RK^G(\underline E G,B) = \varinjlim_{K\subset G \text{ compact }} KK^G(C_0(P_K(G)),B).\]

\subsection{Descent functor}

There exists a natural transformation $KK^G(A,B)\rightarrow KK(A\rtimes_r G, B\rtimes_r G)$ respecting the Kasparov product, which is called the descent functor. The same statement remains true for maximal crossed products. We recall its construction, which was first stated in \cite{LeGall}.\\
 
Let $A$ and $B$ be two $G$-algebras and $(E,\pi,T)\in E^G(A,B)$ be a $K$-cycle. Define :
\begin{itemize}
\item[$\bullet$] $E_G$ to be the completion of $C_c(G,E)$ with respect to the norm $||f||= \sup_{x\in G^{(0)}} \sum_{G^x} |f(g)|^2$,
\item[$\bullet$] $\pi_G$ to be the image of $\pi:A \rightarrow \mathcal L_B(E)$ under the reduced crossed-product functor,
\item[$\bullet$] $T_G$ to be the image of $T\in\mathcal L(E)$ under the reduced crossed-product functor.
\end{itemize}
Then $(E_G,\pi_G,T_G)\in E(A\rtimes_r G,B\rtimes_r G)$ and the map $(E,\pi,T)\mapsto (E_G,\pi_G,T_G) $ induces a homorphism of abelian groups
\[j_G : KK^G(A,B)\rightarrow KK(A\rtimes_r G,B\rtimes_r G)\]
satisfying $j_G(z\otimes_D z')=j_G(z)\otimes_{D\rtimes_r G} j_G(z')$ for any $z\in KK^G(A,D), z'\in KK^G(D,B)$.\\ 

\subsection{The assembly map}

If $Z$ is a proper $G$-compact $G$-space, $Z\rtimes G$ is a proper groupoid by definition, hence there exists a cutoff function $c : Z\rightarrow [0,1]$ such that $\int_{G^{p(z)}} c(zg)d\lambda^x(g) = 1$. The function $g\mapsto c(r(g))^{\frac{1}{2}}c(s(g))^{\frac{1}{2}} $ defines a projection in $C_0(Z)\rtimes_r G$ which we denote by $\mathcal L_Z$. If $Z=P_K(G)$, then $\mathcal L_Z = \mathcal L_K$.

\begin{definition}
The assembly map for $G$ with coefficients in $B$ is defined as the inductive limit of the maps $\mu_{G,B}^{(Z)} : KK^G(C_0(Z),B)\rightarrow K(B\rtimes_r B)$ given by
\[\mu_{G,B}^{(Z)} (z)=[\mathcal L_Z]\otimes_{C_0(Z)\rtimes G} j_G(z),\]
that is $\mu_{G,B} = \varinjlim \mu_{G,B}^{(Z)}$ (one has to check that theses maps respects the inductive systems, which they do).
\end{definition}

%Another definition : with $Ad$...

\subsection{The Baum-Connes conjecture}

\begin{conj}[Baum-Connes conjecture]
Let $G$ be an étale groupoid, and $A$ a $G$-algebra.\\
The Baum-Connes conjecture for $G$ with coefficients in $A$ is the following claim : \\
\center{$\textbf{BC}_{G,A}$ : $\mu_{G,A}$ is an isomorphism of abelian groups.}\\
The Baum-Connes conjecture with coefficients is : for all $G$-algebras $A$, $\mu_{G,A}$ is an isomorphism of abelian groups.
\end{conj}

The conjecture was first stated in \cite{BaumConnes} and in \cite{BaumConnesHigson}. The statement is a generalization of the Connes-Kasparov conjecture, which is simply, when one knows the Baum-Connes conjecture, the latter for almost connected groups (i.e. locally compact groups $G$ such that $G/G_0$ is compact, where $G_0$ is the connected component of the identity).\\

Here is the status of the conjecture :
\begin{itemize}
\item[$\bullet$] The Connes-Kasparov conjecture was established by J. Chabert, S. Echterhoff and R. Nest, who proved that $\mu_{G}$ is an isomorphism for every secound countable almost connected group $G$, and every group of $k$-rational points of a linear algebraic group over a local field of characteristic $0$.
\item[$\bullet$] N. Higson and G. Kasparov proved that the conjecture with coefficients holds for groups having Haagerup property. 
\item[$\bullet$] V. Lafforgue proved the conjecture with coefficients for hyperbolic groups.
\item[$\bullet$] V. Lafforgue showed the conjecture for every semi-simple Lie group, real or $p$-adic reductive, and discrete cocompact subgroup of real rank $1$ Lie groups or of $SL(3,k)$ for any local field $k$. \\
\end{itemize}

Recall that a function $h : G\rightarrow \R$ is said to be of negative type if :
\begin{itemize}
\item[$\bullet$] $h_{|G^{(0)}}=0$,
\item[$\bullet$] $h(g)=h(g^{-1})$ for all $g\in G$,
\item[$\bullet$] for all $t_i\geq 0, i=1,..,n$ such that $\sum t_i=0$ and all $g_1,...,g_n\in G$ having the same range, $\sum_i t_i t_j h(g_i^{-1} g_j )\leq 0$.
\end{itemize}
The function $h$ is said to be locally proper if, for every compact $K \subset G^{(0)}$, $h_{|G_K^K}$ is proper.\\

\begin{definition}
A topological groupoid $G$ is said to be a-T-menable if it satisfies one of the following equivalent property :
\begin{itemize}
\item[$\bullet$] there exists a continuous field of Hilbert spaces over $G^{(0)}$ with a proper affine action of $G$,
\item[$\bullet$] there exists a locally proper negative type function on $G$.
\end{itemize} 
\end{definition}

For nice topological groupoids, the conjecture is known to be false in full generality \cite{HigsonLaffSk}, but the following result was shown by J-L. Tu in \cite{TuThese}
\begin{thm}\label{Tu}
Let $G$ be a locally compact $\sigma$-compact Hausdorff groupoid. Then if $G$ is a-T-menable, the Baum-Connes conjecture with coefficients holds for $G$.
\end{thm}

\subsection{The Coarse Assembly map}

The construction of the Coarse Assembly map is very similar to the Assembly map for groupoids, but relies on different functors.\\

In \cite{SkTuYu} is defined the Roe algebra $C^*(X,B)$ of a discrete metric space with bounded geometry $X$ with coefficients in an arbitrary $C^*$-algebra $B$. For $R>0$, define $C_R[X,B]$ to be the involutive subspace of operators $T\in \mathcal L_B(l^2(X)\otimes H\otimes B)$ such that $T_{xy}\in \mathfrak K(H\otimes B)$ and $T_{xy}=0$ as soon as $d(x,y)>R$, and $\C[X,B]= \cup_{R>0} C_R[X,B]$. 

\begin{definition}
The Roe algebra $C^*(X,B)$ is defined as the Hibert $B$-module obtained after the completion of $\C[X,B]$ with respect to the operator norm in $\mathcal L_B(l^2(X)\otimes H\otimes B)$.
\end{definition} 

This construction is functorial : any $*$-homomorphism $\phi: A\rightarrow B$ gives rise to a $*$-homomorphism $\phi_X : C^*(X,A)\rightarrow C^*(X,B)$. Moreover this functoriality extends to $KK$-theory : there exists a natural transformation 
\[\sigma_X : KK(A,B)\rightarrow KK(C^*(X,A),C^*(X,B))\]
which respects the Kasparov product, i.e. $\sigma(z\otimes_D z')=\sigma_X(z)\otimes_{C^*(X,D)} \sigma_X(z')\quad,\forall z\in KK(A,D),z'\in KK(D,B)$.\\

Let $X$ be a coarse space with bounded geometry. The Rips complex is the inductive system of finite dimensional simplicial complexes
\[P_E(X)=\{\eta\in \text{Prob}(X)\text{ s.t. supp }\eta\subset E\}\quad,\forall E\in \mathcal E_X,\]
endowed with the $*$-weak topology. It is an inductive system with respect to inclusion of entourages.\\
 
The coarse homology group of $X$ with coefficients in an arbitrary $C^*$-algebra $B$ is defined as 
\[KX_*(X,B)=\varinjlim_{E\in \mathcal E_X} RK(P_E(X),B)\]
where the limit is taken along the inductive system $\{P_E(X)\}_{E\in\mathcal E_X}$. \\

Let $E\in \mathcal E_X$. Any simplex $\eta\in P_E(X)$ can be written as a finite sum $\eta=\sum_x \lambda_x(\eta)\delta_x$ by boundedness of the geometry. Here, $\delta_x$ is the Dirac measure at $x\in X$. The functions $\lambda_x :\eta\mapsto \lambda_x(\eta)$ are continuous and satifisfy $\sum\lambda_x(\eta) = 1,\forall\eta\in P_E(X)$, so that
\[V_0\left\{\begin{array}{rcl} 
 C_0(P_E(X)) & \rightarrow 	& l^2(X)\otimes C_0(P_E(X)) 		\\ 
 f           & \mapsto 		& (\lambda_x^{\frac{1}{2}}f)_{x\in X} 
\end{array}\right.\] 
is an isometry of $C_0(P_E(X))$-modules, and $V_0V_0^*\otimes id_H$ defines a projection in $l^2(X)\otimes H\otimes C_0(P_E(X))$ with finite propagation, so class $[\mathcal L_E]\in K_0(C^*(X,C_0(P_E(X)))$.\\

\begin{definition}
The Coarse Assembly map $\mu:KX_*(X,B)\rightarrow K_*(C^*(X,B))$ for $X$ with coefficients in $B$ is defined as 
\[\forall z\in RK(P_E(X), B), \mu_{X,B}(z)=[\mathcal L_E]\otimes_{C^*(X,C_0(P_E(X)))} \sigma_X(z).\]
\end{definition}

The Coarse Baum-Connes conjecture is the following claim.\\

\begin{conj}[Coarse Baum-Connes conjecture]
For any coarse space $X$ with bounded geometry, the Coarse Assembly map $\mu_{X,\C}$ is an isomorphism.\\
\end{conj}

Here are some remarks about the status of the conjecture :
\begin{itemize}
\item[$\bullet$] the conjecture is known to hold for any coarse space that admits a coarse embedding into a  separable Hilbert space\cite{Yu2}.
\item[$\bullet$] couterexamples have been constructed \cite{HigsonLaffSk},
\item[$\bullet$] Let $\Gamma$ be a finitely generated group, and $|\Gamma|$ the coarse space associated to the word-length metric (any two left invariant metric on $\Gamma$ are coarsely equivalent). Then the Coarse Baum-Connes conjecture for $|\Gamma|$ implies the Novikov conjecture for higher signatures for $\Gamma$.\\
\end{itemize}

Here is another description of the Coarse Assembly map which will be of some interest for us.\\

Let $E\in \mathcal E_X$ be an entourage. Take a cycle $(H,\pi, T)\in E(C_0(P_E(X)), \C)$, and equip the finite dimensional simplicial complex $P_E(X)$ with a metric which restricts to the spherical metric on simplices, so that $H$ is a $P_E(X)$-module that we can suppose standard non-degenrate. Then 
\[T'=\sum \lambda_x^{\frac{1}{2}} T \lambda_x^{\frac{1}{2}} \] 
is a compact perturbation of $T$ in $\mathcal L(H)$ with bounded propagation, and is invertible modulo $C^*(P_E(X),H)$, so defines a class $[T']\in K_0(C^*(P_E(X),H))$. But $P_E(X)$ and $X$ are coarsely equivalent, the barycentric map $P_E(X)\rightarrow X$ is a coarse equivalence which is covered by $V_0$, so that 
\[(Ad_{V_0})_* : K_*(C^*(P_E(X),H))\rightarrow K_*(C^*(X,l^2(X)\otimes H))\]
is an isomorphism, which does only depend on the coarse class of the barycentric map.\\

With that in mind, one can show that $\mu_{X}([H,\pi,T])=(Ad_{V_0})_*[T']$.\\

The point of the article \cite{SkTuYu} is to show that the Coarse Assembly map is equivalent to the Assembly map for the coarse groupoid $G=G(X)$ with coefficients in the $G$-algebra $l^\infty=l^\infty(X,\mathfrak K)$. \\

More precisely, $Z=P_{\overline E}(G)$ is a proper cocompact $G$-space with ismomorphic fibers $Z_x\simeq P_E(X)$, and $l^\infty_x \simeq \mathfrak K$, so that the inclusion of groupoid $\{x\}\hookrightarrow G $ induces a morphism $\iota_Z : RK^G(Z,l^\infty)\rightarrow KK(Z_x,\mathfrak K)$ which is actually an isomorphism. Recall that one can construct an $*$-isomorphism $\Psi_X: l^\infty \rtimes_r G\rightarrow C^*(X)$, and $\mu_X^{(Z_x)}\circ \iota_Z = (\Psi_X)_*\circ \mu_{G,l^\infty}^{(Z)}$ holds and respects inductive limits over $Z$ to give the following commutative diagram :
\[\begin{tikzcd}
K_*^{top}(\underline E G, l^\infty) \arrow{r}{\mu_{G,l^\infty}}\arrow{d}{\varinjlim \iota_Z} & K_*(l^\infty\rtimes_r G)\arrow{d}{(\Psi_X)_*} \\
KX_*(X)\arrow{r}{\mu_X} & K_*(C^*(X))
\end{tikzcd}\]  
with vertical arrows being isomorphism.\\

A key point is to use how analytical properties of $G$ translate coarse properties of $X$. In particular, 
\begin{itemize}
\item[$\bullet$] in \cite{SkTuYu} is proved that $X$ admits a coarse embedding into Hilbert space iff $G(X)$ is a-T-menable, 
\item[$\bullet$] in \cite{SkTuYu} is proved that $X$ has property A iff $G(X)$ is amenable, 
\item[$\bullet$] in \cite{FinnSellFibred}, M. Finn-Sell shows that if $X$ admits a coarse fibered embedding, then $G_\partial = G(X)_{|\partial \beta X}$ is a-T-menable. 
\end{itemize}
All of these results, combined with theorem \ref{Tu}, give examples of coarse spaces satisfying the Coarse Baum-Connes conjecture and the Boundary Coarse Baum-Connes conjecture.


















\section{Enoncés quantitatifs}
\begin{frame}
  \tableofcontents[currentsection]
\end{frame}
\begin{frame}{Enoncés quantitatifs}
Deuxième étape : formuler une version controlée de la conjecture de Baum-Connes.\\
\vspace{0.3 cm}
Pour cela, il nous faut d'abord introduire les propriétés suivantes :\\
\begin{itemize}
\item[$\bullet$] $QI_{G,B}(E,E',F,\varepsilon)$ : pour tout $x\in RK^G(P_E(G), B )$, alors $\mu^{\varepsilon,E,F}_{G,B}(x) = 0$ implique que $q_E^{E'}(x)=0$ dans $RK^G(P_{E'}(G),B)$.
\item[$\bullet$] $QS_{G,B}(E,F,F',\varepsilon,\varepsilon')$ : pour tout $y\in K^{\varepsilon,F}(B\rtimes G)$, il existe $x\in RK^G(P_E(G),B)$ tel que $\mu^{\varepsilon',E,F'}_{G,B}(x)=\iota_{\varepsilon,F}^{\varepsilon',F'}(y)$.\\
\end{itemize} 
\end{frame}

\begin{frame}{Enoncés quantitatifs}
Les énoncés quantitatifs forment le résultat central de la thèse. Ils prennent la forme suivante. Soit $G$ un groupoïde étale $\sigma$-compact à base compacte.\\
\vspace{0.3 cm}
\begin{thmfr}[Injectivité Quantitative]
Soient $B$ une $G$-algèbre, et $\tilde B = l^\infty(\N,B\otimes\mathfrak K)$. Alors $\mu_{G,\tilde B}$ est injective ssi pour tout $E\in\mathcal E,\varepsilon\in(0,\frac{1}{4})$ et $F$ tel que $k_J(\varepsilon).E\subseteq F$, il existe $E' \in\mathcal E$ tel que $E\subseteq E'$ et $QI_{G,B}(E,E',\varepsilon,F)$ soit vérifiée.
\end{thmfr}
\end{frame}

\begin{frame}{Enoncés quantitatifs}
\begin{thmfr}[Surjectivité Quantitative]
Soient $B$ une $G$-algèbre, et $\tilde B = l^\infty(\N,B\otimes\mathfrak K)$. Alors il existe $\lambda \geq 1$ tel que $\mu_{G,\tilde B}$ est surjective ssi pour tout $0<\varepsilon<\frac{1}{4\lambda}$ et $F\in\mathcal E$, il existe $E\in\mathcal E$ et $F'$ tel que $k_J .E \subseteq F$ tel que $QS_{B,G}(E,F,F',\varepsilon,\lambda\varepsilon)$ soit vérifiée.
\end{thmfr}
\end{frame}

\begin{frame}{Enoncés quantitatifs}
Ils admettent une version uniforme :
\begin{thmfr}[Version uniforme] Soit $G$ un groupoïde étale $\sigma$-compact avec une base compacte. \\
\begin{itemize}
\item[$\bullet$] Supposons que pour toute $G$-algèbre $B$, $\mu_{G,B}$ soit injective. Alors, pour tout $\varepsilon\in (0,\frac{1}{4})$ et tout $E,F\in\mathcal E$ tels que $k_J(\varepsilon). E\subseteq F$, il existe $E'\in\mathcal E$ tel que $E\subseteq E'$ et tel que $QI_{G,B}(E,E',\varepsilon,F)$ soit satisfait pout toute $G$-algèbre $B$.\\
\item[$\bullet$] Supposons que pour toute $G$-algèbre $B$, $\mu_{G,B}$ soit surjective. Alors, pour un certain $\lambda \geq 1$ et pour tout $\varepsilon\in (0,\frac{1}{4\lambda})$ et tout $F\in\mathcal E$, il existe $E,F'\in\mathcal E$ tels que $k_J(\varepsilon). E\subseteq F'$ et $F\subseteq F'$ tel que, pour toute $G$-algèbre $B$, $QS_{G,B}(E, F,F',\varepsilon,\lambda \varepsilon)$ soit satisfait.
\end{itemize}
\end{thmfr}
\end{frame}




\section{Géométrie Coarse}
\begin{frame}
  \tableofcontents[currentsection]
\end{frame}
\section{Coarse geometry, Roe algebras and étale groupoids}

\subsection{From coarse spaces to $C^*$-algebras}

\begin{definition}
A coarse space is a couple $(X,\mathcal E)$, where $X$ is a set and $\mathcal E$ a coarse structure on the pair groupoid $X\times X$. A coarse map is a map respecting the coarse structure i.e. $h : (X,\mathcal E)\rightarrow (Y,\mathcal F)$ such that $\forall F\in\mathcal F,(h\times h)^{-1}(F)\in \mathcal E$.
\end{definition}

A metric space $(X,d)$ naturally inherits a coarse structure from its bounded subsets :
\[\mathcal E_X = \{E\subset X\times X : \sup_{E} d(x,y)<\infty\}.\]

\begin{definition}
Let $(X,\mathcal E)$ be a coarse space. A $X$-module is a pair $(H_X,\phi )$ where 
\begin{itemize}
\item[$\bullet$] $H_X$ is a Hilbert space, 
\item[$\bullet$] $\phi : C_0(X)\rightarrow \mathcal L(H_X)$ is a $*$-homomorphism.
\item[$\bullet$] The module is said to be non-degenerate if $\{\phi(f)\eta : f\in C_0(X),\eta\in H_X\}$ is dense in $H_X$, and standard if no non-zero function of $C_0X)$ acts as a compact operator of $H_X$.
\end{itemize}
\end{definition}

\textbf{Examples} If $X$ is a discrete metric space with bounded geometry, $l^2(X)\otimes H$ with action by multiplication on the first factor defines a s.n.d. $X$-module.\\
If $X$ is a compact metric space endowed with a measure $\mu$ without atoms, then $L^(X,\mu)$ is a s.n.d. $X$-module.\\

\begin{definition}
Let $T\in \mathcal L(H_X, H_Y)$ be an operator.
\begin{itemize}
\item[$\bullet$] $T$ is said to be locally compact if $\phi(g)T$ and $T\phi(f)$ are compact operators for all $f\in C_0(X),g\in C_0(Y)$.
\item[$\bullet$] The support of $T$ is the complement of the set of points $(x,y)\in X\times X$ such that there exist $f_x,f_y\in C_0(X),C_0(Y)$ such that $f_x(x)\neq 0,f_{y}(y)\neq 0$ and $\phi(f_{y}) T \phi(f_x)=0$.
\item[$\bullet$] The propagation of $T$ is the smallest $E\in \mathcal E$ such that supp $T \subset E$.\textbf{ADAPTER}
\end{itemize}
\end{definition}

Let us define the Roe algebra of $X$ form a fixed s.n.d. $X$-module $H_X$. It will be shown that, up to unnatural isomorphism, it does not depend on the choice of the s.n.d. $X$-module.\\

\begin{definition}
For any $E\in \mathcal E$, define the subspace of locally compact operators with propagation $E$-controlled :
\[C_E[X,H_X] = \{T\in\mathcal L(H_X) : T \text{ is locally compact and supp }T\subset E\}.\]
The Roe algebra is the $C^*$-algebra 
\[C^*(X,H_X) = \overline{\cup_{E\in\mathcal E} C_E[X,H_X]}\]
the closure being taken with respect to the operator norm of $\mathcal L(H_X)$.
\end{definition}

\begin{prop}
Let $(X,\mathcal E),(Y,\mathcal F)$ be two coarse spaces, $H_X,H_Y$ two s.n.d. modules over $X$ and $Y$ repsectively, and $h :X\rightarrow Y$ a coarse map. Then, for any $F\in \mathcal F$, there exists an isometry $V\in \mathcal L(H_X,H_Y)$ such that
\[\text{supp }V \subset (h\times id_Y)^{-1}(E).\]
\end{prop}

\begin{dem}
Extend the representation $\phi$ to $\tilde \phi : L^\infty (X)\rightarrow \mathcal L(H_X)$. \\

Let $F\in \mathcal F$ and $\mathcal U$ be a Borel partition of $Y$ such that $\sqcup_{U\in \mathcal U} U\times U \subset F$ and every $U\in \mathcal U$ is of non-empty interior. If $\chi_A$ denotes the characteristic function of $A$, and because the modules are s.n.d., we can find an isometry \[V^{(U)} :\chi_{h^{-1}(U)}H_X \rightarrow \chi_U H_Y\] 
and by standardness $H_X =\bigoplus \chi(h^{-1}(U)) H_X$ and $H_Y =\bigoplus \chi_U H_Y$, so $V = \oplus V^{(U)}$ fits. Indeed, 
\[\text{supp }V\subset \sqcup h^{-1}(U)\times U \subset (h\times id_Y)^{-1}(F).\] 
\qed
\end{dem}

Now, if $H_X$ and $H'_X$ are two s.n.d. $X$-modules, apply the preceding lemma to the identity map to have an isometry $V: H_X\rightarrow H'_X$ which is supported as close as you want of the diagonal. This induces $Ad_V : C^*(X,H_X)\rightarrow C^*(X,H'_X)$ by $Ad_V(T) = VTV^*$, and gives our isomorphism, which is canonical in $K$-theory. \\

If we decide to fix a s.n.d. module for every coarse space, we can speak of "the" Roe algebra of $X$, and we saw that a coarse map between two coarse spaces $h : X\rightarrow Y$ induces a $*$-homomorphism $h_* : C^*(X,H_X)\rightarrow C^*(Y,H_Y)$. \\

We will now define a dimension on coarse spaces, after an idea of Gromov \textbf{à vérifier}, and that was extensively used in coarse geometry.\\

\begin{definition}
A coarse space $(X,\mathcal E)$ is said to have asymptotic dimension less than $d$ if, for every $E\in \mathcal E$, there is a controlled set $F\in\mathcal F$ and a family $\mathcal U$ of susbsets such that 
\begin{itemize}
\item[$\bullet$] $\mathcal U$ covers $X$,
\item[$\bullet$] every $U\in \mathcal U$ is $F$-controlled, i.e. $\sqcup_{U \in \mathcal U} U\times U \subset F$,
\item[$\bullet$] we have a decomposition $\mathcal U = \mathcal U_0 \sqcup ... \sqcup \mathcal U_d$ such that every pair of subsets of a $\mathcal U^{(j)}$ are $E$-separated.
\end{itemize}
\end{definition}

\subsection{From coarse spaces to groupoids}

This section gives details on he construction of the coarse groupoid from \cite{SkTuYu}.\\

Let $(X,\mathcal E)$ be a coarse space. The idea is to construct an étale groupoid $G(X)$ extending the pair groupoid $X\times X$. As a topological space, 
\[G(X)= \cup_{E\in \mathcal E} \overline{E},\]
the closure being taken in $\beta (X\times X)$, so that it is Hausdorff and locally compact. The base space is $G(X)^{(0)}=\beta X$. Now, the firt and second projections can be extended by universal property of the Stone-Cech compactification to give the source and the range map $s,r : G(X)\rightrightarrows \beta X$ respectively, as shown in this commutative diagramm :  
\[\begin{tikzcd}
X\times X \arrow{r}\arrow{d}{\iota_{X\times X}} &  X \arrow{r}{\iota_X} & \beta X \\
\beta (X\times X) \arrow[dotted]{urr} & & \\ 
\end{tikzcd}.\]

The only non trivial part is to extend the multiplication of the pair groupoid. This is done by extending the inclusion $E \rightarrow X\times X$ to $\overline{E} \rightarrow \beta X\times \beta X$. The corollary 10.31 of \cite{RoeCoarse} assures that 
\begin{lem}
The map $r\times s : \overline E \rightarrow \beta X\times \beta X$ is a topological embedding.
\end{lem}
Using this lemma, we can embed $G(X)$ in $\beta X\times \beta X$ and use the pair multiplication in this groupoid, the point being that such a multiplication $G(X)\times_{s,r} G(X) \rightarrow G(X)$ is continuous and extend the pair multiplication of $X\times X$.\\

The following proposition (Proposition $3.5$ from \cite{SkTuYu}) assures that $X\mapsto G(X)$ is a functor from uniformly locally finite coarse spaces to groupoids with generalized morphisms as arrows.\\

\begin{prop}
Let $(X,\mathcal E_X)$ and $(Y,\mathcal E_Y)$ be two coarse spaces with uniformly locally finite coarse structures, then a coarse map $h :X\rightarrow Y$ induces a generalized morphism from $G(X)$ to $G(Y)$ :
\[\begin{tikzcd}
G(X)\arrow[xshift=-0.7ex]{d}\arrow[xshift=0.7ex]{d} \ & \arrow{dl} Z \arrow{dr} & G(Y)\arrow[xshift=-0.7ex]{d}\arrow[xshift=0.7ex]{d} \\
\beta X & &\beta Y \\
\end{tikzcd}\] 
\end{prop}

We now state a result from \cite{SkTuYu} (lemma $4.4$), and give, for the reader's convenience, a more detailed proof than that of the paper.\\

\begin{prop}
Let $(X,\mathcal E)$ be a uniformly locally finite coarse space, then we have an isomorphism of $C^*$-algebras
\[C^*(X) \simeq l^\infty(X,\mathfrak K) \rtimes_r G(X).\]
\end{prop}

\begin{proof}
Let $D=l^\infty(X,\mathfrak K)$ and $G=G(X)$. The $C^*$-algebra $D\rtimes_r G$ is generated by continuous functions $f : \overline E \rightarrow D$ such that $f(g)\in D_{s(g)}\simeq \mathfrak K$ for some $E\in \mathcal E$. The crossed product is obtained as the closure in the norm operator defined by the actions of such functions by convolution on $\mathcal E = L^2(G,D)$, which defines by definition a faithful map $D\rtimes_r G \rightarrow \mathcal L(\mathcal E)$.\\

Now take the $G$-invariant ideal $J= C_0(X,\mathfrak K)$. As $D\rightarrow \mathcal M(J)$ is faithful, $\mathcal L(\mathcal E)\rightarrow \mathcal L(\mathcal E\otimes_{D} J)$ is isometric, and we obtain a faithful map by composition $D\rtimes_r G \rightarrow \mathcal L(\mathcal E\otimes_D J)$, and $\mathcal E\otimes_D J \simeq L^2(G,J)$.\\

But $C_0(X)$ acts faithfully by multiplication on $C_0(X,\mathfrak K)$, hence on $H_X=L^(G,J)$, which makes it a n.d.s. $X$-module and induces a faithful map $C^*(X,H_X) \rightarrow \mathcal L(H_X)$.\\

If $T\in C_E[X,H_X]$, define $f(g)=T_{s(g),r(g)}$ when $g\in X\times X$, and extend $f$ by continuity to get $f : \overline E \rightarrow D$, and $f(g)\in D_{s(g)}$.\\

If $f:\overline E \rightarrow D$ is a continuous function such that $f(g)\in D_{s(g)}$, define $T\in C_E[X,H_X]$ by $T\xi = f\xi$. This defines a $*$-homomorphism $D\rtimes_r G \rightarrow C^*(X,H_X)$. The previous construction shows it is surjective. But the action of $D\rtimes_r G$ on $H_X$ being faithful, it is also injective, which concludes the proof.\\
\end{proof}














\section{Formule de Künneth}
\begin{frame}
  \tableofcontents[currentsection]
\end{frame}
\section{A quantitative Künneth formula}

In this section, we present an application of the previous results to a Künneth formula for crossed products by an étale groupoid.\\ 

Recall that for $A$ and $B$ two $C^*$-algebras, one can define a homomorphism
\[\alpha_{A,B} : K_*(A)\otimes K_*(B)\rightarrow K_*(A\otimes B) \quad ; \quad (x,y)\mapsto x\otimes   \tau_A(y),\]
where $\tau_A$ is the external Kapsarov product.

\begin{definition}
A $C^*$-algebra $A$ is said to satisfy the Künneth formula if, for every $C^*$-algebra $B$ such that $K_*(B)$ is a free abelian group, $\alpha_{A,B}$ is an isomorphism.
\end{definition}

If $A$ satifies the Künneth formula, then, for any $C^*$-algebra $B$, one has the following exact sequence
\[\begin{tikzcd}[column sep = small] 0 \arrow{r} & K_*(A)\otimes K_*(B)\arrow{r} & K_*(A\otimes B) \arrow{r} & Tor(K_*(A),K_*(B))\arrow{r} & 0 \end{tikzcd}\]

If $(A,\mathcal E)$ is now filtered, $(A\otimes B,\mathcal E)$ is also filtered, and one can define a controlled morphism
\[\hat\alpha_{A,B} : \hat K_*(A)\otimes K_*(B)\rightarrow \hat K_*(A\otimes B) \quad ; \quad (x,y)\mapsto \hat\tau_A(y)(x),\]
which induces $\alpha_{A,B}$ in $K$-theory. \\

\begin{definition}
A filtered $C^*$-algebra $(A,\mathcal E)$ is said to satisfy the quantitative Künneth formula if, for every $C^*$-algebra $B$ such that $K_*(B)$ is a free abelian group, $\hat\alpha_{A,B}$ is a controlled isomorphism.
\end{definition}

The remainder of the section is devoted to prove the following theorem.

\begin{thm}\label{Kunneth}
Let $G$ be an $\sigma$-compact étale groupoid and $A$ a $G$-algebra. Suppose that 
\begin{itemize}
\item[$\bullet$] $G$ satisfies the Baum-Connes conjecture with coefficients,
\item[$\bullet$] for every compact subgroupoid $H$ of $G$, $A\rtimes_r H$ satisfies the Künneth formula.
\end{itemize} 
Then $A\rtimes_r G$ satisfies the quantitative Künneth formula.
\end{thm}

The following corollary is obvious from the theorem.\cite{FinnSellFibred}

\begin{cor}
Let $X$ be a proper discrete metric space with bounded geometry which admits a fibered coarse embedding into a Hilbert space and $A$ a $C^*$-algebra. If $A$ satisfies the Künneth formula, then $A\otimes K(l^2(X))$ satifies the quantitative Künneth formula.
\end{cor}

%%%%%%%%%%%%%%%%%%%%%%%%%%%%%%%%%%%%%%%%%%%%%%%%
\subsection{Induction and Restriction functors}
%%%%%%%%%%%%%%%%%%%%%%%%%%%%%%%%%%%%%%%%%%%%%%%%

We recall the following construction from \cite{LeGall}. Let $G$ and $G'$ two étale groupoids, and $\phi=(Z,p,p')$ a generalized morphism from $G$ to $G'$. \\

For any $G'$-equivariant $B$-Hilbert module $E$, define $\phi^*E=l^2_Z\otimes_{C(G')} E$ where $l^2_Z $ is the completion of the space of sections $C_c(Z,B)$ such that $f(z)\in B_{p'(z)}$. Then $\phi^*E$ is a $G$-equivariant $\phi^*B$-Hilbert module. \\

If $\pi : A\rightarrow \mathcal L_B(E)$ and every $F\in\mathcal L_B(E)$, define $\pi' = \pi\otimes_{C(G')} id_{C_0(G^{(0)})}$, and $F'=F\otimes id_{C_0(G^{(0)})}= $. If $(E,\pi,F)\in E^{G'}(A,B)$, we use the notation $\phi^* (E,\pi,F)=(\phi^*E,\pi',F')$.\\

Le Gall proves \cite{LeGall} that $(E,\pi,F)\mapsto \phi^*(E,\pi,F)$ descends to a homomorphism of abelian groups $KK^{G'}(A,B)\rightarrow KK^{G}(\phi^*A,\phi^*B)$ which respects Kasparov products, i.e. $\phi^*(x\otimes_D y)= \phi^*(x)\otimes_{\phi^*D}\phi^*(y)$ holds for all $x,y\in KK^{G'}(A,B)$.\\

Moreover, the association $(E,\pi)\mapsto (\phi^*E, \pi')$ is functorial at the level of $C^*$-correspondences.

\begin{definition}
$G$ and $G'$ are said to be Morita equivalent when there are two generalized morphisms $Z : G\rightarrow G'$ and $Z' : G'\rightarrow G$ such that their compositions are cohomologous to identities.  
\end{definition}

If $H$ is a subgroupoid of $G$, $G$ naturally carries left and right actions of $G/H\rtimes G$ and $H$, giving rise to two generalized morphisms $\phi : G/H\rtimes G\rightarrow H$ and $\psi : H \rightarrow G/H\rtimes G$, which are inverse of each others ($\phi \psi$ and $\psi\phi$ are cohomologous to the identity). In one sentence, $G/H\rtimes G$ and $H$ are Morita equivalent.\\

If $H$ is compact, one can also see $G$ as a $H$ space, and consider the generalized inclusion  
\[\iota : 
\begin{tikzcd}[column sep = small] 
H \arrow{d} \arrow{d} & \arrow{dl} G \arrow{dr} & G\arrow{d} \arrow{d} \\
 H^{(0)} & & G^{(0)}
\end{tikzcd}\]
and $\eta$ the opposite of $\iota$.  
\begin{definition}
The restriction and induction functors are defined by
\[ \text{Res}_H^G = \iota^*;\quad \text{Ind}_H^G = \eta^*.\]
These are functorial w.r.t. $C^*$-correspondences, and also w.r.t. $KK$-theory.
\end{definition}

The Morita equivalence $G/H\rtimes H \sim H$ allows us to give a very concrete definition of these functors, and actually help to prove that, when $H$ is a compact subgroupoid of $G$ and the action of $G$ is induced by $H$ , $\text{Res}_{H}^G : KK^G(A,B)\rightarrow KK^H( \text{Res}_{H}^G A , \text{Res}_{H}^G B)$ is an isomorphism of abelian groups for any $G$-algebras $A$ and $B$. This is locally satisfied for example when $G$ is ample, i.e. étale with totally disconnected base space.

\begin{definition}
The groupoid $G$ is said to be locally induced if there exists an open cover $\mathcal U$ of $G^{(0)}$ such that for all $U\in\mathcal U$, there exists a compact groupoid $H_U$ acting on $G_{|U}$ and an ismomorphism of groupoids $G_{|U }\simeq G_{|U}/ H_U \rtimes H_U$.  
\end{definition}

Since $G_{|U}/ H_U \rtimes H_U$ is Morita equivalent to $H_U$ and $G$ is Morita equivalent to the pull-back $G[\mathcal U]$, we have an isomorphism 
\[KK^G(A,B) \simeq \prod_{U\in\mathcal U} KK^{H_U}(\phi^* A,\phi^* B).\]

Ample groupoids are example of locally induced groupoid. Recall that an étale groupoid is ample iff the open compact bisections form  a basis of the groupoid. Now take $U$ an open compact bisection, and set $H_U = \{ g\in G\text{ s.t. } gU\cap U \neq \emptyset \}$, which is a compact subgroupoid of $G$ that fulfills the condition.\\

The author presently doesn't kow of any other examples of locally induced groupoids which are not ample. Maybe looking into groupoids with finite asymptotic dimension. 

%%%%%%%%%%%%%%%%%%%%%%%%%%%%%%%%%%%%%%%%%%%%%%%%
\subsection{Induction and Restriction functors}
%%%%%%%%%%%%%%%%%%%%%%%%%%%%%%%%%%%%%%%%%%%%%%%%

\begin{definition} A subset $H\subseteq G$ is called a subgroupoid if :
\begin{itemize}
\item[$\bullet$] for every $x\in G^{(0)}$, $e_x\in H$,
\item[$\bullet$] for all $h,h'\in H $ such that $s(h') = r(h)$, $h'h \in H$,
\item[$\bullet$] if $h\in H$, $h^{-1}\in H$.
\end{itemize}
Then, the restriction of the multiplication, inverse, unit, target and source maps on $H$ defines a structure of groupoid on $H$ over $H^{(0)} = s(H)$. If $G$ is étale, $H$ is also étale. We will write $H< G$ to indicate that $H$ is a subgroupoid of $G$.
\end{definition}

In this section, we define for all subgroupoids $H < G$ induction and restriction transformations. Let $G$ be an étale groupoid and $H<G$.\\

Let $A$ be a $H$-algebra, with action given by $\alpha : s^*A \rightarrow r^* A$. Put $\phi=r\circ \iota$ where $\iota : H^{(0)}\hookrightarrow G^{(0)}$ is the canonical inclusion, hence $\phi^*A = A\otimes_\phi C_0(G)$ is a $C(G)$-algebra . Define the induced $C(G^{(0)})$-algebra :
%\[\text{Ind}_H^G (A) = \{f \in C_0(G,A) \text{ s.t. } h^{-1} f(gh) = f(g),\forall h\in H\}. \]
\[\text{Ind}_H^G (A) = \{f \in \phi^* A \text{ s.t. } \alpha_{h^{-1}}(f_{hg}) = f_g,\forall h\in H,g\in G^{s(h)}\}. \]
Left translation defines an action of $G$ on $\text{Ind}_H^G A$, so that $\text{Ind}_H^G A$ is a $G$-algebra.\\

Let $E$ be a Hilbert $A$-module, endowed with an action $V\in\mathcal L_A(s^*E,r^*E)$ of $H$. Define the induced Hilbert module :
%\[\text{Ind}_H^G (E) = \{f \in C_0(G,E) \text{ s.t. } h^{-1} f(gh) = f(g),\forall h\in H\}. \]
\[\text{Ind}_H^G (E) = \{f \in \phi^* E \text{ s.t. } \alpha_{h^{-1}}(f_{hg}) = f_g,\forall h\in H,g\in G^{s(h)}\}. \]
Left translation defines an action of $G$ on $\text{Ind}_H^G E$, hence $\text{Ind}_H^G E$ is a $G$-equivariant Hilbert $\text{Ind}_H^G A$-module.

Let $A$ and $B$ be two $H$-algebras. For all $(E,\pi,T)\in \mathbb E^G(A,B)$, define $\text{Ind}_H^G \pi$ and $\text{Ind}_H^G T$ as pointwise evaluation and multiplication by $\pi$ and $T$, i.e. $\text{Ind}_H^G T = T\otimes_\phi 1$ and $\text{Ind}_H^G \pi  = \pi\otimes_\phi id$.

\begin{definition}
For all subgroupoids $H<G$, let $A$, $B$ and $D$ be $G$-algebras and $A'$, $B'$ and $D'$ be $H$-algebras. Then, the map $(E,\pi,T)\mapsto ( \text{Ind}_H^G E, \text{Ind}_H^G\pi ,\text{Ind}_H^G T )$ induces an even homomorphism of $\Z_2$-graded abelian groups
\[\text{Ind}_H^G : KK_*^H(A',B')\rightarrow KK_*^G( \text{Ind}_H^G A', \text{Ind}_H^G B') \] 
called the induction transformation.\\
Moreover, by forgetting the action, we naturally have an even homomorphism of $\Z_2$-graded abelian groups 
\[\text{Res}_H^G : KK_*^G(A,B)\rightarrow KK_*^H( \text{Res}_H^G A, \text{Res}_H^G B) \] 
called the restriction transformation by restricting the action.\\
These two transformations respect the Kasparov product, i.e. 
\[ \text{Ind}_H^G(z\otimes z') = \text{Ind}_H^G(z)\otimes_{\text{Ind}_H^G(D')}\text{Ind}_H^G(z')\quad \forall z\in KK^H(A',D'),z'\in KK^H(D',B') \]
and 
\[ \text{Res}_H^G(z\otimes z') = \text{Res}_H^G(z)\otimes_{\text{Res}_H^G(D)}\text{Res}_H^G(z')\quad \forall z\in KK^G(A,D),z'\in KK^G(D,B) \]
The reader is refered to \cite{LeGall} for a proof.
\end{definition}   

Let $Z$ be a right $H$-space. Define on $Z\times_{p,r} G$ the following equivalence relation :
\[(z,g)\sim_H (z.h, h^{-1}g)\quad \forall z\in Z, h\in H^{p(z)},g\in G^{p(z)}.\]

\begin{definition}
The induced $G$-space of a $H$-space $Z$ is defined as $Z\times_H G = (Z\times G) / \sim_H$. 
\end{definition}

Notice that we have a natural identification between $\text{Ind}_H^G C_0(Z)$ and $C_0(Z\times_H G)$.\\

\begin{lem} \label{Restriction} Let $H$ be an open subgroupoid of $G$, and $V$ a $H$-space such that the anchor map $p : V\rightarrow H^{(0)}$ is locally injective. Then, for every $H$-algebra $A$ and every $G$-algebra $B$, the transformations $Res_H^G$ and $Ind_H^G$ induce an isomorphism of $\Z_2$-graded abelian groups :
\[RK^G( G\times_H V, B) \cong RK^H(V, Res_H^G B).\]
\end{lem}

\begin{dem}
It is clear that induction followed by restriction is the identity. For the converse, let $(E,\pi,T)\in \mathbb E^G(C_0(G\times_H V),B)$. The moment map is locally injective, hence, by lemma \ref{JLTform}, we can suppose that $T$ is self-adjoint $G$-equivariant and commutes with the action of $C_0(G\times_H V)$. As $H$ is open in $G$, $x\mapsto (e_{p(x)},x)$ is a topological embedding and $V$ can be seen as a $H$-invariant open subset in $G\times_H V$. Denote by $E_V$ the Hilbert $H$-invariant $B$-submodule of $E$ generated by 
\[\{\pi(f)\xi \ ,f\in C_0(V), \xi\in E\}.\]
Then, $E = \oplus_{g\in G/H} E_{gV}$. Notice that $E_{V}$ is a $H$-equivariant Hilbert $Res_H^G(B)$-module, such that $E\cong Ind_H^G (E_V)$. Moreover, $\pi$ is $G$-equivariant, hence $\pi(a) = Ind_H^G (\pi(a)_{|E_V} )$. As $[T,\pi(a)]=0$ for every $a\in C_0(G\times_H V)$, $T(E_V)\subseteq E_V$, and by $G$-equivariance, $T_{|E_V}$ determines $T$. Hence $T= Ind_H^G (T_{|E_V})$. Hence, if $z=[E,\pi,T]$ and $z_H =[E_V,\pi(a)_{|E_V},T_{|E_V}]$, we proved that $z = Ind_H^G( z_H)$, hence $Ind_H^G \circ Res_H^G (z)= z$.\\
\qed  
\end{dem}

%\begin{dem} 
%Let $A$ be a $H$-algebra, and $B$ a $G$-algebra.
%Let us first notice that the statement holds for equivariant $*$-homomorphisms. Namely,\\

%\begin{itemize} 
%\item[$\bullet$] if $\phi : Ind_H^G(A) \rightarrow B$ is a $G$-equivariant $*$-homomorphism, then $A = A\otimes C_0(G^{(0)})$ is a $H$-invariant subalgebra of $Ind_H^G(A)$, and the restriction $\phi_H$ of $\phi$ to $A$ is a $H$-equivariant $*$-homomorphism ;   
%\item[$\bullet$] if $\psi : A \rightarrow Res_H^G(B)$ is a $H$-equivariant $*$-homomorphism, then $\psi_G = \psi \otimes id_{C_0(G)}$ is a $H$-equivariant $*$-homomorphism.\\
%\end{itemize}

%And these constructions satisfy $(\phi_H)_G = \phi$ and $(\psi_G)_H = \phi$. Moreover, in $KK$-theory, one has $Res_H^G([\phi]) = [\phi_H]\in KK^H(Ind_H^G(A), B)$ and $Ind_H^G([\psi]) = [\psi_H]\in KK^H( A,Res_H^G (B))$.\\

%Let $z\in KK^H(Ind_H^G(A), B)$. As the restriction and induction transformations respect Kasparov products, by property $(d)$ and naturality of the boundary map, we can suppose that $z$ is the inverse in $KK$-theory of a $H$-equivariant $*$-homomorphism. Hence there exists a $H$-equivariant $*$-homomorphism $\phi : B \rightarrow Ind_H^G(A)$ such that $z\otimes_B [\phi] = 1_{Ind_H^G(A)}$ and $[\phi]\otimes_{Ind_H^G(A)} z  = 1_B$. Taking the induction, we get $Ind_H^G(z) = Ind_H^G([\phi])^{-1} = [\phi_G]^{-1}$, and $Res_H^G \circ Ind_H^G(z) = z$. We can similarly prove that $Ind_H^G \circ Res_H^G (z') = z'$ for every $z\in KK^G(A,Res_H^G(B))$.\\
%\qed
%\end{dem}

%\begin{prop} Let $H$ be an open compact subgroupoid of $G$, $U$ a $H$-space and $B$ a $H$-algebra. Then :
%\[Res_H^G : RK^G(G\times_H U , B)\rightarrow RK^H(U , Res_H^G(B)) \]
%is an isomorphism of $\Z_2$-graded abelian groups.
%\end{prop}

\subsection{Strongly proper groupoids}

We now introduce a property on groupoids that will entail a nice result on induction and restriction transformations at the level of $K$-homology.

\begin{definition}\label{StronglyProper}
A groupoid $G$ is said to be strongly proper if there exists an open cover $\mathcal U$ of $G^{(0)}$ such that, for all $U\in\mathcal U$, there exists a compact open subgroupoid $H_U$ of $G$ and a $H_U$-space $Z_U$ together with a $G$-equivariant homeomorphism
\[\psi_U : U \rightarrow G\times_{H_U} Z_U.\] 
An action of $G$ on a space $Z$ is said to be strongly proper if the groupoid $Z\rtimes G$ is strongly proper. A groupoid is said to be in the class $\mathcal C$ if every proper action of $G$ is strongly proper.
\end{definition}

\begin{rk}
For any strongly proper action of $G$ on a space $Z$, there exists an open cover of $Z$ by subsets of the type $V=G\times_H U$, where $H$ is a compact open subgroupoid and $U$ is a $H$-space. Then, by the previous section, we have an isomorphism
\[RK^G(V,B)\cong RK^H(U, Res_H^G (B))\]
for every $G$-algebra $B$. 
\end{rk}
%\begin{prop}
%Let $G$ be a strongly proper groupoid, and $\mathcal U$ an open cover satisfying the conditions of definition \ref{StronglyProper}. Then, the $G$-map $\psi_U $ extends to an isomorphism of groupoids
%\[\begin{tikzcd}
% G_{|U}\arrow{r}\arrow[shift right]{d}\arrow[shift left]{d} & Z_U\rtimes H_U \arrow[shift right]{d}\arrow[shift left]{d} \\
% U \arrow{r}{\psi_U} & G\times_{H_U} Z_U 
%\end{tikzcd}.\]
% Ce morceau bug sous windows : pourquoi ?
%\end{prop}

Let us give examples of groupoids in class $\mathcal C$. Recall the following definition from \cite{Renault} (page $20$).

\begin{definition}
A topological groupoid is said to be ample if it has a basis $G^a$ of neighborhoods consisting of compact open susbets.
\end{definition}

In \cite{paterson} (page $17$) is stated the following property. An étale groupoid $G$ is ample iff $G^{(0)}$ is totally disconnected. Hence the coarse groupoid of every coarse space $X$ is ample, its basis being $\beta X$.

\begin{prop}
Every ample groupoid is in class $\mathcal C$.
\end{prop}

\begin{dem} The following argument was explained to me by Christian Bönicke.\\ 
Let $G$ be an étale ample groupoid and $Z$ a $G$-space with proper action of $G$. Since $G$ is ample, we can cover $Z$ by compact open subsets. Let $\mathcal U$ be such an open cover. For each $U\in \mathcal U$, there exists $V\subseteq Z$ compact open such that $U= G.V$. Put $H= (r\times s)^{-1}(V,V)$, which is, by properness, a compact open subgroupoid of $Z\rtimes G$. Moreover, $U\cong G\times_H p(V)$ $G$-equivariantly. Hence $Z$ is covered by open subsets of the form $G\times_H V$, with $H$ being compact open subgroupoids and $V$ being $H$-spaces.\\ 
\qed
\end{dem}

\begin{rk}
Let $\Gamma$ be a discrete group. Then every proper action of $\Gamma$ on a space $Z$ is strongly proper by definition.
\end{rk}


%The author presently doesn't know of any other examples of strongly proper groupoids which are not ample.% Maybe looking into groupoids with finite asymptotic dimension. 

%Ample groupoids are example of locally induced groupoid. Recall that an étale groupoid is ample iff the open compact bisections form  a basis of the groupoid. Now take $U$ an open compact bisection, and set $H_U = \{ g\in G\text{ s.t. } gU\cap U \neq \emptyset \}$, which is a compact subgroupoid of $G$ that fulfills the condition.\\

%\begin{definition}
%The groupoid $G$ is said to be locally induced if there exists an open cover $\mathcal U$ of $G^{(0)}$ such that for all $U\in\mathcal U$, there exists a compact groupoid $H_U$ acting on $G_{|U}$ and an ismomorphism of groupoids $G_{|U }\simeq G_{|U}/ H_U \rtimes H_U$.  
%\end{definition}

%Since $G_{|U}/ H_U \rtimes H_U$ is Morita equivalent to $H_U$ and $G$ is Morita equivalent to the pull-back $G[\mathcal U]$, we have an isomorphism 
%\[KK^G(A,B) \simeq \prod_{U\in\mathcal U} KK^{H_U}(\phi^* A,\phi^* B).\]






%%%%%%%%%%%%%%%%%%%%%%%%%%%%%%%%%%%%%%%%%%%%%%%%%%%%
\subsection{Baum-Connes and the Künneth formula}
%%%%%%%%%%%%%%%%%%%%%%%%%%%%%%%%%%%%%%%%%%%%%%%%%%%%
The first step in proving theorem \ref{Kunneth} is to define a analytical version of $\alpha_{A,B}$ when an étale groupoid $G$ is given. More precisely, we first construct a homomorphism $\alpha_{A,B}^G : K_*^{top}(G,A)\otimes K_*(B)\rightarrow K_*^{top}(G,A\otimes B )$, inductive limit of $\alpha_{A,B}^{G,Z} : RK^G(Z,A)\otimes K_*(B)\rightarrow RK^G(Z,A\otimes B )$ where $Z$ runs through $G$-proper $G$-compact spaces. Then we show that the controlled assembly map intertwines $\alpha^{G,P_d(G)}_{A,B}$ and $\alpha_{A\rtimes_r G,B}$ uniformly.\\

If $A$ is a $G$-algebra, and $B$ a $C^*$-algebra, $A\otimes B$ naturally inherits a $G$-algebra structure with trivial action of $G$ on the $B$ factor.\\

Let $Z$ be a $G$-proper space. Define the homomorphism :
\[\alpha_{A,B}^{G,Z} : RK^G_*(Z,A)\otimes K_*(B)\rightarrow RK_*^G(Z,A\otimes B) \quad ; \quad (x,y)\mapsto x\otimes_{}   \tau_A(y),\]
which respects inductive limits w.r.t. inclusions of $G$-proper spaces, so that it induces
\[\alpha_{A,B}^G : K_*^{top}(G,A)\otimes K_*(B)\rightarrow K_*^{top}(G,A\otimes B ).\]

To prove theorem \ref{Kunneth}, we will need the following result.

\begin{thm}\label{TopologicalKunneth}
Let $E$ be an open relatively compact subset of $G$ and $P_E$. If for all compact subgroupoids $H$ of $G$, $A\rtimes_r H$ satisfies the Künneth formula, then $\alpha_{A,B}^{G,P_E(G)}$ is an isomorphism for all $C^*$-algebras $B$ such that $K_*(B)$ is a free abelian group.
\end{thm}

\begin{proof}
Let $Z_0\subset Z_1\subset ... \subset Z_n $ be the skeleton decomposition of $P_E(G)$.\\

Let us prove by induction that $\alpha_j=\alpha^{G,Z_j}_{A,B}$ is an isomorphism. By a standard argument similar to the proof of theorem \ref{prod}, it is sufficient to prove the statement for $j=0$.\\

Let $U\subset Z_0$ be a $G$-compact $G$-proper space. There exists a compact subgroupoid $H$ of $G$ and a proper $H$-space such that $U = \text{Ind}_H^G V$. The following diagram is commutative :
\[\begin{tikzcd}[column sep = small] RK_*^G(U,A)\otimes K_*(B) \arrow{r}{\alpha_{A,B}^{G,U}} \arrow{d}{\text{Res}_H^G} & RK^G(U, A\otimes B) \arrow{d}{\text{Res}_H^G}\\
RK_*^H(V, \text{Res}_H^GA)\otimes K_*(B) \arrow{r}{\alpha_{\text{Res}_H^G A,B}^{H,V}} \arrow{d}{\mu_{H,\text{Res}_H^G A \otimes id}} & 
	RK_*^H(V, \text{Res}_H^G A\otimes B)\arrow{d}{\mu_{H,\text{Res}_H^G A }} \\
K_*(\text{Res}_H^G A\rtimes_r H)\otimes K_*(B)\arrow{r}{\alpha_{\text{Res}_H^G A,B}} & K_*\left((\text{Res}_H^G A\otimes B)\rtimes_r H\right) 
\end{tikzcd}\]
The first vertical arrows are isomorphisms by property of the restriction functor, the last horizontal one and the last vertical ones are by hypothesis, so that $\alpha_{A,B}^{G,U}$ is an isomorphism. By taking the inductive limit on $G$-proper $G$-compact spaces $U\subset Z_0$, we get that $\alpha_{A,B}^{G,Z_0}$ is an isomorphism.
\end{proof}

\begin{lem}\label{KunnethLemma}
There exists a positive number $\kappa$ and a function $c : (0,\frac{1}{4\kappa})\times\mathcal E\rightarrow \mathcal E$, decreasing in $E$ and non increasing in $\varepsilon$ such that, for any étale groupoid $G$, any filtered $G$-algebra $(A,\mathcal E)$, and any $C^*$-algebra $B$, the following diagram commutes :
\[\begin{tikzcd}
RK_*^G(P_E(G),A)\otimes K_*(B) \arrow{r}{\alpha^{G,P_E(G)}_{A,B}}\arrow{d}{\mu_{G,A}^{\varepsilon,F}\otimes \text{id}_{K_*(B)}} & 
RK_*^G(P_E(G),A\otimes B) \arrow{d}{\mu_{G,A}^{\alpha_ \tau\varepsilon,h_\varepsilon F}} \\
K_*^{\varepsilon,F}(A\rtimes_r G)\otimes K_*(B) \arrow{r}{\alpha_{A\rtimes_r G,B}^{\varepsilon,F}} & 
K_*^{\alpha_\tau\varepsilon,h_{\varepsilon}F}((A\otimes B)\rtimes_r G)\\
\end{tikzcd}\] 
for all $\varepsilon\in(0,\frac{1}{4\kappa})$ and $E,F\in\mathcal E$ such that $F\geq c(\varepsilon,E)$. 
\end{lem}

\begin{dem}
Let $z\in RK_*^G(P_E(G),A)$ and $y\in K_*(B)$. As the action of $G$ on $B$ is trivial, $(A\rtimes_r G)\otimes B\simeq (A\otimes B)\rtimes_r G$ so that $\hat\tau_A(y)\circ \hat J_{G}(z) \sim \hat J_{G}(z)\circ \hat\tau_A(y) $ holds. This entails
\[\hat J_G^{\alpha_\tau\varepsilon, h_{\tau,\varepsilon}E}(z)\left([\mathcal L_E,0]_{\varepsilon, E} \right)\otimes \tau_{A\rtimes_r G}(y) = 
\left( \hat\tau(y)\circ \hat J_G(z) \right)^{\varepsilon, E}\left([\mathcal L_E,0]_{\varepsilon, E}\right),\]
which is just the statement of the lemma.\\
\qed
\end{dem}

%\textbf{SHOW THAT $\hat J_G(z)\circ \hat\tau(y)=$}

\begin{dem} \textit{of theorem \ref{Kunneth}}\\
Let $\kappa>0$ and $c : (0,\frac{1}{\kappa})\times\mathcal E\rightarrow \mathcal E$ as in lemma \ref{KunnethLemma}.

Fix two $G$-algebras $A$ and a $C^*$-algebra $B$, seen as a $G$-algebra with trivial action. As $(A\otimes B)\rtimes_r G \simeq (A\rtimes_r G)\otimes B$, we will make the slight abuse of coufounding them. Let us prove that $\hat \alpha_{A\rtimes_r G,B}$ is a controlled isomorphism.\\

Let $F\in\mathcal E$ and $\varepsilon\in (0,\frac{1}{4})$. Choose $E \in\mathcal E$ such that $F\geq R(\varepsilon,E)$. By controlled surjectivity, there exists $F'\in \mathcal E$ such that $QS_{G,A\otimes B}(\varepsilon,\lambda\varepsilon, F,F',E)$ holds.\\

Let $y\in K^{\varepsilon,F}_*((A\otimes B)\rtimes_r G)$. By controlled surjectivity of $\hat\mu_{G,A\otimes B}$, there exists $z\in RK_*^G(P_E(G), A\otimes B) $ such that $\mu_{G,A\otimes B}^{\kappa\varepsilon,F'}(z)=\iota_{\varepsilon,F}^{\kappa\varepsilon,F'}(y)$. Theorem \ref{TopologicalKunneth} entails that there exists $z_1 \in RK_*^G(P_E(G),A)\otimes K_*(B) $ such that $\alpha_{A,B}^{G,P_E(G)}(z_1)=z$, and lemma \ref{KunnethLemma} shows that, if $z_2=(\mu_{G,A}^{\varepsilon,E,F} \otimes id_{\hat K_*(B)})(z_1)$, then $\alpha_{A\rtimes_r G,B}(z_2)=\iota_{}^{}(y)$.\\

Let us show the controlled injectivity. Let $\varepsilon\in(0,\frac{1}{4})$ and $F\in\mathcal E$.\\ 

Let $x\in K^{\varepsilon,F}_*(A\rtimes_r G)\otimes K_*(B)$ such that $\alpha_{A\rtimes_r G,B}^{\varepsilon,F}(x)=0$ in $K^{\alpha_\tau,h_{\tau,\varepsilon} F}((A\otimes B)\rtimes_r G)$. By controlled surjectivity of $\hat\mu_{G,A}$, there exists $z\in RK(P_E(G),A)\otimes K_*(B)$ such that $\mu_{G,A}^{\alpha_\tau,h_{\tau,\alpha_\tau\varepsilon} F}(z)=(\iota_{\varepsilon,F}^{\varepsilon',F'}\otimes id) (x)$. If $z' = \alpha_{A,B}^{G,P_E(G)}(z)$, lemma \ref{KunnethLemma} entails that $\mu^{E}_{G,A\otimes B}(z)=0$, which by controlled injectivity implies that $q_E^{E'}(z')=0$ in $RK^G(P_{E'}(G),A\otimes B)$. As $\alpha_{A,B}^{G,P_E(G)}$ is an isomorphism, $q_E^{E'}(z)=0$ in $RK^G(P_{E'}(G),A)\otimes K(B)$.

\end{dem}






























\footnotesize
\begin{frame}{Références}
\bibliographystyle{plain}
\bibliography{biblio2}
\end{frame} 

\Large
\begin{frame}{}
Merci de votre attention !
\end{frame}

\end{document}
